

\setcounter{chapter}{8}

\chapter{多项式环}

我们以对前两章事实的总结开始这一章关于多项式环的讨论（必要时给出引用）。基本定义在 7.2 节中给出得稍详细一些。为方便起见，环 \(R\) 总设为带单位元 \(1 \neq 0\) 的交换环。

\section{定义与基本性质}

系数取自 \(R\) 的未定元 \(x\) 的\textbf{多项式环} \(R[x]\) 是所有形式和 \(a_nx^n + a_{n-1}x^{n-1} + \cdots + a_1x + a_0\)（\(n \geq 0\) 且每个 \(a_i \in R\)）的集合。若 \(a_n \neq 0\)，则多项式次数为 \(n\)，\(a_nx^n\) 是首项，\(a_n\) 是首项系数（零多项式的首项系数定义为 0）。若 \(a_n = 1\)，则多项式是首一的。多项式的加法是“逐分量”的：
\[
\sum_{i=0}^{n} a_ix^i + \sum_{i=0}^{n} b_ix^i = \sum_{i=0}^{n}(a_i+b_i)x^i
\]
（这里为了定义不同次数的多项式的加法，\(a_n\) 或 \(b_n\) 可以为零）。乘法首先定义 \((ax^i)(bx^j) = abx^{i+j}\)，然后通过分配律推广到所有多项式，故一般有
\[
\left(\sum_{i=0}^{n} a_ix^i\right)\left(\sum_{j=0}^{m} b_jx^j\right) = \sum_{k=0}^{n+m}\left(\sum_{i+j=k} a_ib_j\right)x^k.
\]
这样 \(R[x]\) 是带单位元（来自 \(R\) 的单位元）的交换环，其中我们把 \(R\) 与常多项式的子环等同。

我们已经注意到，若 \(R\) 是整环，则两个多项式之积的首项是各因子首项之积。下面是 7.2 节的命题 4，我们为完整起见在此记录。

\begin{proposition}\label{pro:9.1}
设 \(R\) 是整环。则

（1）若 \(p(x), q(x)\) 非零，则 \(\deg p(x)q(x) = \deg p(x) + \deg q(x)\)；

（2）\(R[x]\) 的单位就是 \(R\) 的单位；

（3）\(R[x]\) 是整环。
\end{proposition}

还要回忆若 \(R\) 是整环，则 \(R[x]\) 的商域由所有商 \(\frac{p(x)}{q(x)}\) 组成，其中 \(q(x)\) 不是零多项式（称为系数在 \(R\) 中的、\(x\) 的有理函数域）。

下一个结果描述 \(R\) 的理想与 \(R[x]\) 的理想之间的关系。

\begin{proposition}\label{pro:9.2}
设 \(I\) 是环 \(R\) 的理想，\((I) = I[x]\) 表示由 \(I\) 生成的 \(R[x]\) 的理想（即系数在 \(I\) 中的多项式集合）。则
\[
R[x]/(I) \cong (R/I)[x].
\]
特别地，若 \(I\) 是 \(R\) 的素理想，则 \((I)\) 是 \(R[x]\) 的素理想。
\end{proposition}

\begin{proof}
存在自然的映射 \(\varphi : R[x] \to (R/I)[x]\)，由把多项式的每个系数模 \(I\) 约化给出。这两个环中加法与乘法的定义表明 \(\varphi\) 是环同态。核恰好是每个系数都是 \(I\) 的元素的多项式集合，即 \(\ker\varphi = I[x] = (I)\)，这证明了命题的第一部分。最后一个陈述由命题 \ref{pro:9.1} 推出，因为若 \(I\) 是 \(R\) 中的素理想，则 \(R/I\) 是整环，故 \((R/I)[x]\) 也是整环。这表明若 \(I\) 是 \(R\) 的素理想，则 \((I)\) 是 \(R[x]\) 的素理想。
\end{proof}

注意：若 \(I\) 是 \(R\) 的极大理想，则 \((I)\) 不必是 \(R[x]\) 的极大理想。然而若 \(I\) 在 \(R\) 中极大，则由 \(I\) 和 \(x\) 生成的 \(R[x]\) 的理想在 \(R[x]\) 中极大。

我们现在给出命题 \ref{pro:9.2} 的“约化同态”的一个例子，它在后面许多场合有用（“约化同态”也在 7.3 节末就模 \(n\) 约化整数讨论过）。

\begin{example}
设 \(R = \mathbb{Z}\)，考虑 \(\mathbb{Z}\) 的理想 \(n\mathbb{Z}\). 则上面的同构可以写成
\[
\mathbb{Z}[x]/n\mathbb{Z}[x] \cong \mathbb{Z}/n\mathbb{Z}[x],
\]
而由把系数模 \(n\) 约化给出的自然投影映射 \(\mathbb{Z}[x] \to \mathbb{Z}/n\mathbb{Z}[x]\) 是环同态。若 \(n\) 是合数，则商环不是整环。然而若 \(n\) 是素数 \(p\)，则 \(\mathbb{Z}/p\mathbb{Z}\) 是域，故 \(\mathbb{Z}/p\mathbb{Z}[x]\) 是整环（事实上是欧几里得整环，如我们不久将看到）。我们还看到系数被 \(p\) 整除的多项式集合是 \(\mathbb{Z}[x]\) 中的素理想。
\end{example}

我们以描述向多变量多项式环的自然推广来结束本节。

\begin{definition}
系数在 \(R\) 中的、变量 \(x_1, x_2, \ldots, x_n\) 的\textbf{多项式环}，记作 \(R[x_1, x_2, \ldots, x_n]\)，由
\[
R[x_1, x_2, \ldots, x_n] = R[x_1, x_2, \ldots, x_{n-1}][x_n]
\]
归纳定义。
\end{definition}

这个定义意味着我们可以把系数在 \(R\) 中的 \(n\) 个变量的多项式简单地看作一个变量（比如 \(x_n\)）的多项式，但其系数本身是 \(n-1\) 个变量的多项式。用一种稍具体的形式表述，系数在 \(R\) 中的、\(x_1, x_2, \ldots, x_n\) 的非零多项式是非零单项式项的有限和，即形如
\[
a x_1^{d_1}x_2^{d_2}\cdots x_n^{d_n}
\]
的元素（\(a \in R\) 是该项的系数，\(d_i\) 是非负整数）的有限和。首一单项式 \(x_1^{d_1}x_2^{d_2}\cdots x_n^{d_n}\) 简单地称为\textbf{单项式}，它是项 \(ax_1^{d_1}\cdots x_n^{d_n}\) 的单项式部分。指数 \(d_i\) 称为该项关于 \(x_i\) 的\textbf{次数}，而和 \(d = d_1+d_2+\cdots+d_n\) 称为该项的\textbf{次数}。有序 \(n\) 元组 \((d_1, d_2, \ldots, d_n)\) 是该项的\textbf{多重次数}。非零多项式的次数是其各单项式项次数的最大值。若一个多项式的所有项次数相同，则称它是\textbf{齐次的}或是\textbf{形式}。若 \(f\) 是 \(n\) 个变量的非零多项式，则 \(f\) 中所有 \(k\) 次单项式项之和称为 \(f\) 的 \(k\) 次\textbf{齐次分量}。若 \(f\) 的次数为 \(d\)，则 \(f\) 可以唯一地写成 \(f_0 + f_1 + \cdots + f_d\)，其中 \(f_k\) 是 \(f\) 的 \(k\) 次齐次分量（\(0 \leq k \leq d\)，某些 \(f_k\) 可以为零）。

最后，为定义系数在 \(R\) 中的任意多个变量的多项式环，我们取上述类型的单项式项（但变量不限于 \(x_1, \ldots, x_n\)）的有限和，配以自然的加法与乘法。另一种方式是把这个环定义为所考虑的有限多个变量上的所有多项式环的并。

\begin{example}
整系数两变量 \(x, y\) 的多项式环 \(\mathbb{Z}[x, y]\) 由形如 \(ax^iy^j\)（次数为 \(i+j\)）的单项式项的所有有限和组成。例如
\[
p(x, y) = 2x^3 + xy - 1 \quad \text{和} \quad q(x, y) = -3xy + 2y^2 + x^2y^3
\]
都是 \(\mathbb{Z}[x, y]\) 的元素，次数分别为 3 和 5. 我们有
\[
p(x, y) + q(x, y) = 2x^3 - 2xy + 2y^2 + x^2y^3 - 1
\]
及
\[
p(x, y)q(x, y) = 2x^5y^3 + x^3y^4 - 6x^4y + 4x^3y^2 - x^2y^3 - 3x^2y^2 + 2xy^3 + 3xy - 2y^2,
\]
一个次数为 8 的多项式。为把这个后一个多项式（比如）看作 \(y\) 的、系数在 \(\mathbb{Z}[x]\) 中的多项式（如上面多变量多项式环的定义所述），我们把多项式写成
\[
(x^3)y^4 + (2x^5 - x^2 + 2x)y^3 + (4x^3 - 3x^2 - 2)y^2 + (-6x^4 + 3x)y.
\]
\(f = f(x, y) = p(x, y)q(x, y)\) 的非零齐次分量是：\(f_2 = 3xy - 2y^2\)（次数 2）、\(f_4 = -3x^2y^2 + 2xy^3\)（次数 4）、\(f_5 = -6x^4y + 4x^3y^2 - x^2y^3\)（次数 5）、\(f_7 = x^3y^4\)（次数 7）和 \(f_8 = 2x^5y^3\)（次数 8）。
\end{example}

命题 \ref{pro:9.1} 中的每个陈述对任意多个变量的多项式环都成立。这可以由有限多个变量的情形归纳推出，而任意多个变量的情形则由用并给出的定义推出。

\subsection*{习题}

\begin{enumerate}
\item 设 \(p(x, y, z) = 2x^2y - 3xy^3z + 4y^2z^5\)、\(q(x, y, z) = 7x^2 + 5x^2y^3z^4 - 3x^2z^3\) 是 \(\mathbb{Z}[x, y, z]\) 中的多项式。

（a）把 \(p\) 和 \(q\) 各写成 \(x\) 的、系数在 \(\mathbb{Z}[y, z]\) 中的多项式。

（b）求 \(p\) 和 \(q\) 的次数。

（c）求 \(p\) 和 \(q\) 关于三个变量 \(x, y, z\) 中每一个的次数。

（d）计算 \(pq\)，并求 \(pq\) 关于三个变量 \(x, y, z\) 中每一个的次数。

（e）把 \(pq\) 写成变量 \(z\) 的、系数在 \(\mathbb{Z}[x, y]\) 中的多项式。

\item 假设 \(p\) 和 \(q\) 的系数在 \(\mathbb{Z}/3\mathbb{Z}\) 中，重复前一习题。

\item 若 \(R\) 是交换环且 \(x_1, x_2, \ldots, x_n\) 是 \(R\) 上相互独立的变量，证明对 \(\{1, 2, \ldots, n\}\) 的任意置换 \(\sigma\)，\(R[x_{\sigma(1)}, \ldots, x_{\sigma(n)}]\) 同构于 \(R[x_1, \ldots, x_n]\).

\item 证明理想 \((x)\) 和 \((x, y)\) 是 \(\mathbb{Q}[x, y]\) 中的素理想，但只有后一个理想是极大理想。

\item 证明 \((x, y)\) 和 \((2, x, y)\) 是 \(\mathbb{Z}[x, y]\) 中的素理想，但只有后一个理想是极大理想。

\item 证明 \((x, y)\) 不是 \(\mathbb{Q}[x, y]\) 中的主理想。

\item 设 \(R\) 是带 1 的交换环。证明 \(R\) 上的多于一个变量的多项式环不是主理想整环。

\item 设 \(F\) 是域，\(R = F[x, x^2y, x^3y^2, \ldots, x^ny^{n-1}, \ldots]\) 是多项式环 \(F[x, y]\) 的子环。

（a）证明 \(R\) 与 \(F[x, y]\) 的分式域相同。

（b）证明 \(R\) 含有一个不是有限生成的理想。

\item 证明系数取自任意交换环的无穷多个变量的多项式环含有不是有限生成的理想。

\item 证明环 \(\mathbb{Z}[x_1, x_2, x_3, \ldots]/(x_1x_2, x_3x_4, x_5x_6, \ldots)\) 含有无穷多个极小素理想（参见第 7.4 节习题 36）。

\item 证明 \(\mathbb{Q}[x, y]\) 中理想 \(I = (x, y^2)\) 的根是 \((x, y)\)（参见第 7.4 节习题 30）。由此推出 \(I\) 是准素理想但不是素理想的幂（参见第 7.4 节习题 41）。

\item 设 \(R = \mathbb{Q}[x, y, z]\)，横线表示过渡到 \(\mathbb{Q}[x, y, z]/(xy - z^2)\). 证明 \(P = (\bar{x}, \bar{z})\) 是素理想。证明 \(\bar{x}\bar{y} \in P^2\)，但 \(P^2\) 中不含 \(\bar{y}\) 的任何幂。（这表明 \(P\) 是素理想，但其平方不是准素理想——参见第 7.4 节习题 41。）

\item 证明对任意域 \(F\)，环 \(F[x, y]/(y^2 - x)\) 与 \(F[x, y]/(y^2 - x^2)\) 不同构。

\item 设 \(R\) 是整环，\(i, j\) 是互素整数。证明理想 \((x^i - y^j)\) 是 \(R[x, y]\) 中的素理想。〔考虑由把 \(x\) 映到 \(t^j\)、把 \(y\) 映到 \(t^i\) 定义的从 \(R[x, y]\) 到 \(R[t]\) 的环同态 \(\varphi\). 证明 \(R[x, y]\) 的元素与 \((x^i - y^j)\) 中某个元素之差是一个关于 \(y\) 的次数至多为 \(j-1\) 的多项式 \(f(x)\)，并注意对 \(0 \leq s < j\)，\(\varphi(x^r y^s)\) 的指数互不相同。〕

\item 设 \(p(x_1, \ldots, x_n)\) 是 \(R[x_1, \ldots, x_n]\) 中的 \(k\) 次齐次多项式。证明对所有 \(\lambda \in R\) 有
\[
p(\lambda x_1, \lambda x_2, \ldots, \lambda x_n) = \lambda^k p(x_1, x_2, \ldots, x_n).
\]

\item 证明两个齐次多项式之积仍是齐次的。

\item \(R[x_1, \ldots, x_n]\) 中的理想 \(I\) 称为\textbf{齐次理想}，若只要 \(p \in I\) 则 \(p\) 的每个齐次分量也在 \(I\) 中。证明一个理想是齐次理想当且仅当它可以由齐次多项式生成。〔用对次数的归纳证明“当”这一方向。〕

\noindent 下面的习题表明在非交换环 \(R\) 上处理多项式时必须小心（\(R[x]\) 中的环运算对非交换环 \(R\) 与对交换环 \(R\) 的定义方式相同），特别是在把多项式看作函数时。

\item 设 \(R\) 是任意环，\(\operatorname{Func}(R)\) 是所有从 \(R\) 到自身的函数构成的环。若 \(p(x) \in R[x]\) 是多项式，设 \(\tilde{p} \in \operatorname{Func}(R)\) 是由 \(\tilde{p}(r) = p(r)\) 定义的 \(R\) 上的函数（即把 \(R[x]\) 中的多项式看作通过“在 \(r\) 处求值”定义 \(R\) 上函数的通常方式）。

（a）对固定的 \(a \in R\)，证明“在 \(a\) 处求值”是从 \(\operatorname{Func}(R)\) 到 \(R\) 的环同态（参见第 7.3 节定理 7 后的例 4）。

（b）证明由 \(\varphi(p(x)) = \tilde{p}\) 定义的映射 \(\varphi : R[x] \to \operatorname{Func}(R)\) 一般不是环同态。由此推出当把多项式看作函数时，多项式恒等式不必给出相应的恒等式。〔若 \(R = \mathbb{H}\) 是实 Hamilton 四元数环，证明 \(p(x) = x^2+1\) 分解为 \((x+i)(x-i)\)，但 \(p(j) = 0\) 而 \((j+i)(j-i) \neq 0\).〕

（c）对固定的 \(a \in R\)，证明（a）和（b）中的映射复合得到的从 \(R[x]\) 到 \(R\) 的“在 \(a\) 处求值”是环同态当且仅当 \(a\) 在 \(R\) 的中心中。
\end{enumerate}

\section{域上的多项式环（一）}

我们现在更仔细地考察系数环是域 \(F\) 的情形。可以通过定义 \(N(p(x)) = p(x)\) 的次数（并令 \(N(0) = 0\)）在 \(F[x]\) 上定义一个范数。由初等代数我们知道，可以用一个（比如）有理系数多项式去除另一个非零有理系数多项式，得到商与余式。对任意域也一样。

\begin{theorem}\label{thm:9.3}
设 \(F\) 是域。则多项式环 \(F[x]\) 是欧几里得整环。更具体地说，若 \(a(x)\) 和 \(b(x)\) 是 \(F[x]\) 中两个多项式且 \(b(x)\) 非零，则存在唯一的 \(q(x), r(x) \in F[x]\) 使得
\[
a(x) = q(x)b(x) + r(x), \quad \text{其中 } r(x) = 0 \text{ 或 } \deg r(x) < \deg b(x).
\]
\end{theorem}

\begin{proof}
若 \(a(x)\) 是零多项式，则取 \(q(x) = r(x) = 0\). 因此我们可以假设 \(a(x) \neq 0\)，并对 \(n = \deg a(x)\) 归纳证明 \(q(x)\) 与 \(r(x)\) 的存在性。设 \(b(x)\) 的次数为 \(m\). 若 \(n < m\)，则取 \(q(x) = 0\)、\(r(x) = a(x)\). 否则 \(n \geq m\). 写
\[
a(x) = a_nx^n + a_{n-1}x^{n-1} + \cdots + a_1x + a_0
\]
和
\[
b(x) = b_mx^m + b_{m-1}x^{m-1} + \cdots + b_1x + b_0.
\]
则多项式 \(a'(x) = a(x) - \frac{a_n}{b_m}x^{n-m}b(x)\) 的次数小于 \(n\)（我们已经安排从 \(a(x)\) 中减去它的首项）。注意这个多项式是有定义的，因为系数取自域且 \(b_m \neq 0\). 由归纳假设，存在多项式 \(q'(x)\) 和 \(r(x)\) 使得
\[
a'(x) = q'(x)b(x) + r(x), \quad \text{其中 } r(x) = 0 \text{ 或 } \deg r(x) < \deg b(x).
\]
于是令 \(q(x) = q'(x) + \frac{a_n}{b_m}x^{n-m}\)，我们得到
\[
a(x) = q(x)b(x) + r(x), \quad \text{其中 } r(x) = 0 \text{ 或 } \deg r(x) < \deg b(x),
\]
这就完成了归纳步骤。

至于唯一性，假设 \(q_1(x)\) 和 \(r_1(x)\) 也满足定理的条件。则 \(a(x) - q(x)b(x)\) 与 \(a(x) - q_1(x)b(x)\) 的次数都小于 \(m = \deg b(x)\). 这两个多项式之差，即 \(b(x)(q(x) - q_1(x))\)，次数也小于 \(m\). 但两个非零多项式之积的次数等于它们次数之和（因为 \(F\) 是整环），故必有 \(q(x) - q_1(x) = 0\)，即 \(q(x) = q_1(x)\). 由此 \(r(x) = r_1(x)\)，完成证明。
\end{proof}

\begin{corollary}\label{cor:9.4}
若 \(F\) 是域，则 \(F[x]\) 是主理想整环，也是唯一分解整环。
\end{corollary}

\begin{proof}
这由上一章的结果立即推出。
\end{proof}

还要回忆第 8.2 节推论 8：若 \(R\) 是任意交换环使得 \(R[x]\) 是主理想整环（或欧几里得整环），则 \(R\) 必是域。然而我们将在下一节看到，只要 \(R\) 本身是唯一分解整环，\(R[x]\) 就是唯一分解整环。

\begin{example}
（1）由上述注记，环 \(\mathbb{Z}[x]\) 不是主理想整环。正如我们已经看到的（7.4 节开头的例 3），理想 \((2, x)\) 在这个环中不是主理想。

（2）\(\mathbb{Q}[x]\) 是主理想整环，因为系数落在域 \(\mathbb{Q}\) 中。由 2 和 \(x\) 在 \(\mathbb{Z}[x]\) 中生成的理想在 \(\mathbb{Q}[x]\) 的子环 \(\mathbb{Z}[x]\) 中不是主理想。然而这个理想在 \(\mathbb{Q}[x]\) 中生成的理想是主的；事实上它是整个环（故以 1 为生成元），因为 2 是 \(\mathbb{Q}[x]\) 中的单位。

（3）若 \(p\) 是素数，则把 \(\mathbb{Z}[x]\) 模素理想 \((p)\) 约化得到的环 \(\mathbb{Z}/p\mathbb{Z}[x]\) 是主理想整环，因为系数落在域 \(\mathbb{Z}/p\mathbb{Z}\) 中。这个例子表明一个不是主理想整环的环的商环可以是主理想整环。为了在这个例子中追踪上面的理想 \((2, x)\)，注意若 \(p = 2\)，则理想 \((2, x)\) 约化为商环 \(\mathbb{Z}/2\mathbb{Z}[x]\) 中的理想 \((x)\)，它是一个真（极大）理想。若 \(p \neq 2\)，则 2 在商环中是单位，故理想 \((2, x)\) 约化为整个环 \(\mathbb{Z}/p\mathbb{Z}[x]\).

（4）\(\mathbb{Q}[x, y]\)，即系数为有理数的二元多项式环，不是主理想整环，因为这个环是 \(\mathbb{Q}[x][y]\) 而 \(\mathbb{Q}[x]\) 不是域（任何正次数元素都不可逆）。理想 \((x, y)\) 在这个环中不是主理想，这留作习题。我们很快会看到 \(\mathbb{Q}[x, y]\) 是唯一分解整环。
\end{example}

我们注意到，除法算法中作用于 \(a(x), b(x) \in F[x]\) 的商与余式在下述意义下与域的扩张无关。设域 \(F\) 包含于域 \(E\) 中，且对某个满足定理 \ref{thm:9.3} 条件的 \(Q(x), R(x)\) 在 \(E[x]\) 中有 \(a(x) = Q(x)b(x) + R(x)\). 写 \(a(x) = q(x)b(x) + r(x)\)，其中 \(q(x), r(x) \in F[x]\)，并在环 \(E[x]\) 中应用定理 \ref{thm:9.3} 的唯一性条件，推出 \(Q(x) = q(x)\) 且 \(R(x) = r(x)\). 特别地，\(b(x)\) 在环 \(E[x]\) 中整除 \(a(x)\) 当且仅当 \(b(x)\) 在 \(F[x]\) 中整除 \(a(x)\). 另外，\(a(x)\) 与 \(b(x)\) 的最大公因子（可由欧几里得算法得到），在通过指定其为首一来使其唯一之后，无论把这些元素看作 \(F[x]\) 中还是 \(E[x]\) 中的元素都是相同的。

\subsection*{习题}

\noindent 设 \(F\) 是域，\(x\) 是 \(F\) 上的未定元。

\begin{enumerate}
\item 设 \(f(x) \in F[x]\) 是次数 \(n \geq 1\) 的多项式，横线表示过渡到商环 \(F[x]/(f(x))\). 证明对每个 \(g(x)\) 都存在唯一的次数 \(\leq n-1\) 的多项式 \(g_0(x)\) 使得 \(\overline{g(x)} = \overline{g_0(x)}\)（等价地，元素 \(1, x, \ldots, x^{n-1}\) 是 \(F\) 上向量空间 \(F[x]/(f(x))\) 的一组基——特别地，这个空间的维数为 \(n\)）。〔利用除法算法。〕

\item 设 \(F\) 是阶为 \(q\) 的有限域，\(f(x)\) 是 \(F[x]\) 中次数 \(n \geq 1\) 的多项式。证明 \(F[x]/(f(x))\) 有 \(q^n\) 个元素。〔利用前一习题。〕

\item 设 \(f(x)\) 是 \(F[x]\) 中的多项式。证明 \(F[x]/(f(x))\) 是域当且仅当 \(f(x)\) 不可约。〔利用第 8.2 节命题 7。〕

\item 设 \(F\) 是有限域。证明 \(F[x]\) 含有无穷多个素元。（注意在无穷域上，次数为 1 的多项式构成多项式环中无穷多个素元。）

\item 列出环 \(F[x]/(p(x))\) 的所有理想，其中 \(F\) 是域、\(p(x)\) 是 \(F[x]\) 中的多项式（用 \(p(x)\) 的分解来描述它们）。

\item 简略描述下列环的环结构：

（a）\(\mathbb{Z}[x]/(2)\)，（b）\(\mathbb{Z}[x]/(x)\)，（c）\(\mathbb{Z}[x]/(x^2)\)，（d）\(\mathbb{Z}[x, y]/(x^2, y^2, 2)\).

证明在最后一个环中对每个 \(a\) 都有 \(a^2 = 0\) 或 1，并确定那些满足 \(a^2 = 0\) 的元素。确定这些环中每一个的特征（参见第 7.3 节习题 26）。

\item 确定环 \(\mathbb{Z}[x]/(2, x^3+1)\) 的所有理想。

\item 求 \(a(x) = x^3-2\) 与 \(b(x) = x+1\) 在 \(\mathbb{Q}[x]\) 中的最大公因子，并把它写成 \(a(x)\) 与 \(b(x)\) 的（\(\mathbb{Q}[x]\) 中的）线性组合。

\item 求 \(a(x) = x^5+2x^3+x^2+x+1\) 与多项式 \(b(x) = x^5+x^4+2x^3+2x^2+2x+1\) 在 \(\mathbb{Q}[x]\) 中的最大公因子，并把它写成 \(a(x)\) 与 \(b(x)\) 的（\(\mathbb{Q}[x]\) 中的）线性组合。

\item 求 \(a(x) = x^3+4x^2+x-6\) 与 \(b(x) = x^5-6x+5\) 在 \(\mathbb{Q}[x]\) 中的最大公因子，并把它写成 \(a(x)\) 与 \(b(x)\) 的（\(\mathbb{Q}[x]\) 中的）线性组合。

\item 设 \(f(x)\) 和 \(g(x)\) 是 \(\mathbb{Q}[x]\) 中两个非零多项式，最大公因子为 \(d(x)\).

（a）给定 \(h(x) \in \mathbb{Q}[x]\)，证明存在多项式 \(a(x), b(x) \in \mathbb{Q}[x]\) 满足方程 \(a(x)f(x) + b(x)g(x) = h(x)\) 当且仅当 \(h(x)\) 被 \(d(x)\) 整除。

（b）若 \(a_0(x), b_0(x) \in \mathbb{Q}[x]\) 是（a）中方程的一组特解，证明该方程的全部解由
\[
a(x) = a_0(x) + m(x)\frac{g(x)}{d(x)}, \qquad b(x) = b_0(x) - m(x)\frac{f(x)}{d(x)}
\]
给出，其中 \(m(x)\) 取遍 \(\mathbb{Q}[x]\) 中的多项式。〔参见第 8.1 节习题 4。〕

\item 设 \(F[x, y_1, y_2, \ldots]\) 是无穷多个变量 \(x, y_1, y_2, \ldots\) 在域 \(F\) 上的多项式环，并设 \(I\) 是这个环中的理想 \((x-y_1^2, y_1-y_2^2, \ldots, y_i-y_{i+1}^2, \ldots)\). 定义 \(R\) 为环 \(F[x, y_1, y_2, \ldots]/I\)，从而在 \(R\) 中每个 \(y_{i+1}\) 的平方是 \(y_i\)，且模 \(I\) 有 \(y_1^2 = x\)，即对每个 \(i\)，\(x\) 都有一个 \(2^i\) 次根。把 \(y_i\) 在 \(R\) 中的像记作 \(x^{1/2^i}\). 设 \(R_n\) 是由 \(F\) 和 \(x^{1/2^n}\) 生成的 \(R\) 的子环。

（a）证明 \(R_1 \subseteq R_2 \subseteq \cdots\)，且 \(R\) 是所有这些 \(R_n\) 的并，即 \(R = \bigcup_{n \geq 1} R_n\).

（b）证明 \(R_n\) 同构于一个变量的多项式环，故 \(R_n\) 是主理想整环。由此推出 \(R\) 是 Bézout 整环（参见第 8.2 节习题 7）。〔先证明环 \(S_n = F[x, y_1, \ldots, y_n]/(x-y_1^2, y_1-y_2^2, \ldots, y_{n-1}-y_n^2)\) 同构于多项式环 \(F[y_n]\). 然后证明 \(y_n\) 在 \(R_n\) 中满足的任何多项式关系都给出某个 \(N \geq n\) 的 \(S_N\) 中的一个相应关系。〕

（c）证明由 \(x, x^{1/2}, x^{1/4}, \ldots\) 在 \(R\) 中生成的理想不是有限生成的（故 \(R\) 不是主理想整环）。

\item 本习题引入一个非交换环，它是“右”欧几里得整环（也是“左”主理想整环），但不是“左”欧几里得整环（也不是“右”主理想整环）。设 \(F\) 是特征为 \(p\) 的域，其中并非每个元素都是 \(p\) 次幂：\(F \neq F^p\)（例如系数在 \(\mathbb{F}_p\) 中的变量 \(t\) 的有理函数域 \(F = \mathbb{F}_p(t)\) 就是这样的域）。设 \(R = F\{x\}\) 是“扭”多项式环，其元素是 \(x\) 的多项式 \(\sum_{i=0}^n a_ix^i\)（系数在 \(F\) 中），加法为通常的（逐项）加法
\[
\sum_{i=0}^n a_ix^i + \sum_{i=0}^n b_ix^i = \sum_{i=0}^n (a_i+b_i)x^i,
\]
但乘法由非交换的法则
\[
\left(\sum_{i=0}^n a_ix^i\right)\left(\sum_{j=0}^m b_jx^j\right) = \sum_{k=0}^{n+m}\left(\sum_{i+j=k} a_ib_j^{p^i}\right)x^k
\]
定义。这个乘法来自定义 \(xa = a^px\)（对每个 \(a \in F\)）（于是 \(x\) 的幂与系数不交换），并以自然方式推广。设 \(N\) 是由取环中多项式的次数定义的范数：\(N(f) = \deg(f)\).

（a）证明对每个 \(a \in F\) 和每个整数 \(k \geq 0\) 有 \(x^ka = a^{p^k}x^k\)，并且在这个乘法定义下 \(R\) 是环。〔利用 \(F\) 的特征为 \(p\) 时对每个 \(a, b \in F\) 有 \((a+b)^p = a^p+b^p\)，故对每个 \(a, b \in F\) 也有 \((a+b)^{p^k} = a^{p^k}+b^{p^k}\).〕

（b）证明 \(R\) 中两个元素之积的次数等于这两个元素次数之和。证明 \(R\) 没有零因子。

（c）证明 \(R\) 关于 \(N\) 是“右欧几里得”的，即对任意 \(f, g \in R\)，\(g \neq 0\)，存在 \(q, r \in R\) 使得
\[
f = qg + r, \quad \text{其中 } r = 0 \text{ 或 } \deg(r) < \deg(g).
\]
由此证明 \(R\) 的每个左理想都是主理想。

（d）设 \(f = ex\)，其中 \(e \in F\)，\(e \notin F^p\)，并设 \(g = x\). 证明不存在 \(q, r \in R\) 使得
\[
f = gq + r, \quad \text{其中 } r = 0 \text{ 或 } \deg(r) < \deg(g),
\]
所以特别地 \(R\) 关于 \(N\) 不是“左欧几里得”的。证明 \(R\) 中由 \(x\) 和 \(ex\) 生成的右理想不是主理想。由此断定 \(R\) 关于任何范数都不是“左欧几里得”的。
\end{enumerate}

\section{唯一分解整环上的多项式环}

我们在命题 \ref{pro:9.1} 中已经看到，若 \(R\) 是整环，则 \(R[x]\) 也是整环。另外，这样的 \(R\) 可以嵌入它的分式域 \(F\)（第 7.5 节定理 15），于是 \(R[x] \subseteq F[x]\) 是子环，而 \(F[x]\) 是欧几里得整环（因而是主理想整环和唯一分解整环）。许多关于 \(R[x]\) 的计算可以在 \(F[x]\) 中完成，代价是允许出现分数系数。这就直接引出如下问题：\(F[x]\) 中的计算（比如多项式的分解）如何用来给出 \(R[x]\) 中的信息？

例如，假设 \(p(x)\) 是 \(R[x]\) 中的多项式。由于 \(F[x]\) 是唯一分解整环，我们可以把 \(p(x)\) 唯一地分解为 \(F[x]\) 中不可约元素之积。自然要问：我们是否能在 \(R[x]\) 中做同样的事，即 \(R[x]\) 是唯一分解整环吗？一般来说答案是否定的，因为若 \(R[x]\) 是唯一分解整环，则常数多项式必须唯一地分解为 \(R[x]\) 中不可约元素之积，由次数可加它们必是 0 次的，也就是说 \(R\) 本身必须是唯一分解整环。于是若 \(R\) 是不是唯一分解整环的整环，则 \(R[x]\) 不可能是唯一分解整环。另一方面，结果表明若 \(R\) 是唯一分解整环，则 \(R[x]\) 也是唯一分解整环。证明的方法是先在 \(F[x]\) 中唯一分解，然后“清分母”得到 \(R[x]\) 中的唯一分解。要使这一点精确，第一步是把 \(F[x]\) 中多项式的分解与 \(R[x]\) 中的分解作比较。

\begin{proposition}\label{pro:9.5}
（\textbf{高斯引理}）设 \(R\) 是唯一分解整环，分式域为 \(F\)，\(p(x) \in R[x]\). 若 \(p(x)\) 在 \(F[x]\) 中可约，则 \(p(x)\) 在 \(R[x]\) 中可约。更精确地说，若对某些非常数多项式 \(A(x), B(x) \in F[x]\) 有 \(p(x) = A(x)B(x)\)，则存在非零元素 \(r, s \in F\) 使得 \(rA(x) = a(x)\) 与 \(sB(x) = b(x)\) 都落在 \(R[x]\) 中，且 \(p(x) = a(x)b(x)\) 是 \(R[x]\) 中的一个分解。
\end{proposition}

\begin{proof}
方程 \(p(x) = A(x)B(x)\) 右端各多项式的系数是域 \(F\) 中的元素，从而是唯一分解整环 \(R\) 中元素的商。把所有这些系数乘上一个公共分母，我们得到方程 \(dp(x) = a'(x)b'(x)\)，其中现在 \(a'(x)\) 和 \(b'(x)\) 都是 \(R[x]\) 的元素，\(d\) 是 \(R\) 中的非零元素。若 \(d\) 是 \(R\) 中的单位，则命题成立，取 \(a(x) = d^{-1}a'(x)\)、\(b(x) = b'(x)\). 假设 \(d\) 不是单位，并把 \(d\) 写成 \(R\) 中不可约元素之积，例如 \(d = p_1\cdots p_n\). 由于 \(p_1\) 在 \(R\) 中不可约，理想 \((p_1)\) 是素理想（参见第 8.3 节命题 12），故由上面的命题 \ref{pro:9.2}，理想 \(p_1R[x]\) 是 \(R[x]\) 中的素理想，且 \((R/p_1R)[x]\) 是整环。把方程 \(dp(x) = a'(x)b'(x)\) 模 \(p_1\) 约化，我们在这个整环中得到方程 \(0 = \bar{a}'(x)\bar{b}'(x)\)（横线表示这些多项式在商环中的像），故两个因子中必有一个为 0，比如说 \(a'(x)\). 但这意味着 \(a'(x)\) 的所有系数都被 \(p_1\) 整除，从而 \(\frac{1}{p_1}a'(x)\) 的系数也在 \(R\) 中。换句话说，在方程 \(dp(x) = a'(x)b'(x)\) 中我们可以从 \(d\)（左端）以及 \(a'(x)\) 或 \(b'(x)\)（右端）之一中消去一个因子 \(p_1\)，而仍得到 \(R[x]\) 中的方程。但现在左端的因子 \(d\) 少了一个不可约因子。对 \(d\) 的其余因子重复同样的做法，我们就能把 \(d\) 的所有因子都消入右端的两个多项式，剩下方程 \(p(x) = a(x)b(x)\)，其中 \(a(x), b(x) \in R[x]\)，且 \(a(x), b(x)\) 分别是 \(A(x), B(x)\) 的 \(F\)-倍数。这完成了证明。
\end{proof}

注意我们不能证明 \(a(x)\) 和 \(b(x)\) 必定分别是 \(A(x)\) 和 \(B(x)\) 的 \(R\)-倍数，因为例如我们可以在 \(\mathbb{Q}[x]\) 中取 \(A(x) = 2x\)、\(B(x) = \frac{x}{2}\) 来把 \(x^2\) 分解，而 \(A(x)\) 和 \(B(x)\) 的任何整数倍都不给出 \(x^2\) 在 \(\mathbb{Z}[x]\) 中的分解。

环 \(R\) 的元素在唯一分解整环 \(F[x]\) 中成为单位（\(F[x]\) 的单位就是 \(F\) 的非零元素）。例如 \(7x\) 在 \(\mathbb{Z}[x]\) 中分解为两个不可约元素之积：7 和 \(x\)（故 \(7x\) 在 \(\mathbb{Z}[x]\) 中不是不可约元素），而 \(7x\) 在 \(\mathbb{Q}[x]\) 中是单位 7 乘以不可约元素 \(x\)（故 \(7x\) 在 \(\mathbb{Q}[x]\) 中不可约）。下面的推论表明，这实质上就是 \(R[x]\) 中不可约元素与 \(F[x]\) 中不可约元素之间的唯一差别。

\begin{corollary}\label{cor:9.6}
设 \(R\) 是唯一分解整环，\(F\) 是它的分式域，\(p(x) \in R[x]\). 假设 \(p(x)\) 各系数的最大公因子是 1. 则 \(p(x)\) 在 \(R[x]\) 中不可约当且仅当它在 \(F[x]\) 中不可约。特别地，若 \(p(x)\) 是首一多项式且在 \(R[x]\) 中不可约，则 \(p(x)\) 在 \(F[x]\) 中不可约。
\end{corollary}

\begin{proof}
由上面的高斯引理，若 \(p(x)\) 在 \(F[x]\) 中可约，则它在 \(R[x]\) 中可约。反之，关于 \(p(x)\) 系数的最大公因子的假设蕴涵：若它在 \(R[x]\) 中可约，则 \(p(x) = a(x)b(x)\)，其中 \(a(x)\) 和 \(b(x)\) 都不是 \(R[x]\) 中的常数多项式。同一个分解说明 \(p(x)\) 在 \(F[x]\) 中可约，完成证明。
\end{proof}

\begin{theorem}\label{thm:9.7}
\(R\) 是唯一分解整环当且仅当 \(R[x]\) 是唯一分解整环。
\end{theorem}

\begin{proof}
上面我们已经指出，\(R[x]\) 是唯一分解整环迫使 \(R\) 是唯一分解整环。反之，假设 \(R\) 是唯一分解整环，\(F\) 是它的分式域，\(p(x)\) 是 \(R[x]\) 的非零元素。设 \(d\) 是 \(p(x)\) 各系数的最大公因子，于是 \(p(x) = dp'(x)\)，其中 \(p'(x)\) 各系数的最大公因子是 1. \(p(x)\) 的这种分解在 \(d\) 的变化（即 \(R\) 中单位的变化）意义下唯一，且由于 \(d\) 可以唯一地分解为 \(R\) 中的不可约元素（它们在更大的环 \(R[x]\) 中也是不可约元素），只需证明 \(p'(x)\) 可以唯一地分解为 \(R[x]\) 中的不可约元素。于是我们可以假设 \(p(x)\) 各系数的最大公因子是 1. 我们还可以假设 \(p(x)\) 不是 \(R[x]\) 中的单位，即 \(\deg p(x) > 0\).

由于 \(F[x]\) 是唯一分解整环，\(p(x)\) 可以唯一地分解为 \(F[x]\) 中不可约元素之积。由高斯引理，这样的分解蕴涵 \(p(x)\) 在 \(R[x]\) 中有一个分解，其因子分别是对应的 \(F[x]\) 中因子的 \(F\)-倍数。由于 \(p(x)\) 各系数的最大公因子是 1，这些 \(R[x]\) 中每个因子各系数的最大公因子必为 1. 由推论 \ref{cor:9.6}，这些因子中每一个都是 \(R[x]\) 中的不可约元素。这表明 \(p(x)\) 可以写成 \(R[x]\) 中不可约元素的有限乘积。

\(p(x)\) 分解的唯一性由 \(F[x]\) 中的唯一性推出。假设
\[
p(x) = q_1(x)\cdots q_r(x) = q'_1(x)\cdots q'_s(x)
\]
是 \(p(x)\) 在 \(R[x]\) 中分解为不可约元素的两种分解。由于 \(p(x)\) 各系数的最大公因子是 1，上面每个不可约因子也具有同样的性质——特别地，每个都有正次数。由推论 \ref{cor:9.6}，每个 \(q_i(x)\) 和 \(q'_j(x)\) 都是 \(F[x]\) 中的不可约元素。由 \(F[x]\) 中的唯一分解，\(r = s\)，且（必要时重新排列后）对每个 \(i \in \{1, \ldots, r\}\)，\(q_i(x)\) 与 \(q'_i(x)\) 在 \(F[x]\) 中相伴。剩下要证明它们在 \(R[x]\) 中相伴。由于 \(F[x]\) 的单位正好是 \(F\) 的非零元素，我们需要考察何时有 \(q(x) = \frac{a}{b}q'(x)\)，其中 \(q(x), q'(x) \in R[x]\)，而 \(a, b\) 是 \(R\) 的非零元素，且 \(q(x)\) 与 \(q'(x)\) 各系数的最大公因子都是 1. 此时 \(bq(x) = aq'(x)\)；左端各系数的最大公因子是 \(b\)，右端是 \(a\). 由于在唯一分解整环中非零多项式各系数的最大公因子在相伴意义下唯一，存在 \(R\) 中的单位 \(u\) 使得 \(a = ub\). 于是 \(q(x) = uq'(x)\)，故 \(q(x)\) 与 \(q'(x)\) 也在 \(R\) 中相伴。这完成了证明。
\end{proof}

\begin{corollary}\label{cor:9.8}
若 \(R\) 是唯一分解整环，则系数在 \(R\) 中的任意多个变量的多项式环也是唯一分解整环。
\end{corollary}

\begin{proof}
对有限多个变量，这由定理 \ref{thm:9.7} 归纳推出，因为 \(n\) 个变量的多项式环可以看作一个变量的、系数在 \(n-1\) 个变量的多项式环中的多项式环。一般情形由“任意多个变量的多项式环是有限多个变量的多项式环之并”这一（上面的）定义推出。
\end{proof}

\begin{example}
（1）\(\mathbb{Z}[x]\)、\(\mathbb{Z}[x, y]\) 等等都是唯一分解整环。环 \(\mathbb{Z}[x]\) 给出一个不是主理想整环的唯一分解整环的例子。

（2）类似地，\(\mathbb{Q}[x]\)、\(\mathbb{Q}[x, y]\) 等等都是唯一分解整环。
\end{example}

我们先前看到，若 \(R\) 是唯一分解整环、分式域为 \(F\)、\(p(x) \in R[x]\)，则我们可以提出 \(p(x)\) 各系数的最大公因子 \(d\)，得到 \(p(x) = dp'(x)\)，其中 \(p'(x)\) 在 \(R[x]\) 和 \(F[x]\) 中都不可约。现在假设 \(R\) 是任意整环，分式域为 \(F\). 在 \(R\) 中最大公因子的概念可能没有意义，然而人们仍可以问：比如说，一个在 \(R[x]\) 中不可约的首一多项式是否仍在 \(F[x]\) 中不可约（即推论 \ref{cor:9.6} 的最后一个断言是否成立）。

首先注意，若首一多项式 \(p(x)\) 可约，则它必在 \(R[x]\) 中有分解 \(p(x) = a(x)b(x)\)，其中 \(a(x)\) 和 \(b(x)\) 都是首一的非常数多项式（回忆 \(p(x)\) 的首项是各因子首项之积，故 \(a(x)\) 和 \(b(x)\) 的首项系数都是单位——于是我们可以把它安排成 1）。换句话说，非常数首一多项式 \(p(x)\) 不可约当且仅当它不能写成两个次数更小的首一多项式之积。

现在我们看到，若 \(R\) 是任意整环、\(p(x)\) 是 \(R[x]\) 中的首一不可约多项式，则 \(p(x)\) 未必在 \(F[x]\) 中不可约。例如，设 \(R = \mathbb{Z}[2i] = \{a+2bi \mid a, b \in \mathbb{Z}\}\)（复数的子环），\(p(x) = x^2+1\). 则 \(R\) 的分式域是 \(F = \{a+bi \mid a, b \in \mathbb{Q}\}\). 多项式 \(p(x)\) 在 \(F[x]\) 中唯一地分解为两个线性因子之积：\(x^2+1 = (x-i)(x+i)\)，所以特别地 \(p(x)\) 在 \(F[x]\) 中可约。这两个因子都不在 \(R[x]\) 中（因为 \(i \notin R\)），故 \(p(x)\) 在 \(R[x]\) 中不可约。特别地，由推论 \ref{cor:9.6}，\(\mathbb{Z}[2i]\) 不是唯一分解整环。

\subsection*{习题}

\begin{enumerate}
\item 设 \(R\) 是整环，商域为 \(F\)，\(p(x)\) 是 \(R[x]\) 中的首一多项式。假设 \(p(x) = a(x)b(x)\)，其中 \(a(x)\) 和 \(b(x)\) 是 \(F[x]\) 中次数小于 \(p(x)\) 的首一多项式。证明若 \(a(x) \notin R[x]\) 则 \(R\) 不是唯一分解整环。由此推出 \(\mathbb{Z}[2\sqrt{2}]\) 不是 U.F.D.

\item 证明若 \(f(x)\) 和 \(g(x)\) 是有理系数多项式，且乘积 \(f(x)g(x)\) 有整数系数，则 \(g(x)\) 的任意系数与 \(f(x)\) 的任意系数之积是整数。

\item 设 \(F\) 是域。证明 \(F[x]\) 中 \(x\) 的系数等于 0 的多项式集合 \(R\) 是 \(F[x]\) 的子环，且 \(R\) 不是 U.F.D.〔证明 \(x^6 = (x^2)^3 = (x^3)^2\) 给出 \(x^6\) 的两个不同不可约分解。〕

\item 设 \(R = \mathbb{Z} + x\mathbb{Q}[x] \subset \mathbb{Q}[x]\) 是常项为整数的有理系数多项式集合。

（a）证明 \(R\) 是整环，其单位是 \(\pm 1\).

（b）证明 \(R\) 中的不可约元素是形如 \(\pm p\)（\(p\) 是 \(\mathbb{Z}\) 中的素数）的元素，以及在 \(\mathbb{Q}[x]\) 中不可约、常项为 \(\pm 1\) 的多项式 \(f(x)\). 证明这些不可约元素在 \(R\) 中是素元。

（c）证明 \(x\) 不能写成 \(R\) 中不可约元素之积（特别地，\(x\) 不是不可约元素），并断定 \(R\) 不是 U.F.D.

（d）证明 \(x\) 在 \(R\) 中不是素元，并描述商环 \(R/(x)\).

\item 设 \(R = \mathbb{Z} + x\mathbb{Q}[x] \subset \mathbb{Q}[x]\) 是上一习题中考虑的环。

（a）假设 \(f(x), g(x) \in \mathbb{Q}[x]\) 是两个非零的有理系数多项式，且 \(x^r\) 是在 \(\mathbb{Q}[x]\) 中同时整除 \(f(x)\) 与 \(g(x)\) 的最高次幂（即 \(r\) 是出现在 \(f(x)\) 或 \(g(x)\) 中的最低次项的次数的较小者）。设 \(f_r\) 与 \(g_r\) 分别是 \(f(x)\) 与 \(g(x)\) 中 \(x^r\) 的系数（由 \(r\) 的定义，其中至少有一个非零）。于是存在整数 \(u, v\) 与某个非零的 \(d_r \in \mathbb{Q}\) 使得 \(uf_r + vg_r = d_r\)（参见第 2.4 节习题 14）。证明存在 \(\mathbb{Q}[x]\) 中的多项式 \(d(x)\)，它是 \(f(x)\) 与 \(g(x)\) 在 \(\mathbb{Q}[x]\) 中的最大公因子，且其最低次项恰为 \(d_rx^r\).

（b）证明 \(f(x) = d(x)q_1(x)\) 与 \(g(x) = d(x)q_2(x)\)，其中 \(q_1(x)\) 和 \(q_2(x)\) 都是 \(\mathbb{Q}[x]\) 的子环 \(R\) 的元素。

（c）证明对 \(R\) 中的多项式 \(a(x), b(x)\) 有 \(d(x) = a(x)f(x) + b(x)g(x)\). 〔在欧几里得整环 \(\mathbb{Q}[x]\) 中 \(a(x), b(x)\) 的存在性立即得到。用第 9.2 节习题 11 证明可以选取 \(a(x)\) 和 \(b(x)\) 使其落在 \(R\) 中。〕

（d）由（a）和（b）推出在 \(\mathbb{Q}[x]\) 中 \(Rf(x) + Rg(x) = Rd(x)\)，并用它证明 \(R\) 是 Bézout 整环（参见第 8.2 节习题 7）。

（e）证明（d）、上一习题的结果以及第 8.3 节习题 11 蕴涵 \(R\) 必含有不是主理想的理想（因而不是有限生成的理想）。事实上证明 \(I = x\mathbb{Q}[x]\) 是 \(R\) 的一个不是有限生成的理想。
\end{enumerate}

\section{不可约性判别法}

若 \(R\) 是唯一分解整环，则由推论 \ref{cor:9.8}，系数在 \(R\) 中的任意多个变量的多项式环也是唯一分解整环。于是确定这样的多项式环中的不可约元素就很有意义，特别是环 \(R[x]\) 中的。在一个变量的情形，非常数首一多项式在 \(R[x]\) 中不可约当且仅当它不能写成两个次数更小的多项式之积。判断一个多项式是否有因子往往很难检验，特别是对多变量的高次多项式。不可约性判别法的目的就是为判断某些类型的多项式是否不可约提供更方便的手段。

我们大部分时候把注意力限制在一个变量的多项式，其系数环是唯一分解整环。由高斯引理，只需考虑 \(F[x]\) 中的分解，其中 \(F\) 是 \(R\) 的分式域（不过我们偶尔也会在系数环只是整环时考虑不可约性问题）。下一个命题考虑何时有一个一次因子（线性因子）。

\begin{proposition}\label{pro:9.9}
设 \(F\) 是域，\(p(x) \in F[x]\). 则 \(p(x)\) 有一个一次因子当且仅当 \(p(x)\) 在 \(F\) 中有根，即存在 \(a \in F\) 使 \(p(a) = 0\).
\end{proposition}

\begin{proof}
若 \(p(x)\) 有一个一次因子，则由于 \(F\) 是域，我们可以设这个因子是首一的，即形如 \((x-a)\)（\(a \in F\)）的。但此时 \(p(a) = 0\).

反之，假设 \(p(a) = 0\). 由 \(F[x]\) 中的除法算法，我们可以写
\[
p(x) = q(x)(x-a) + r,
\]
其中 \(r\) 是常数。由于 \(p(a) = 0\)，\(r\) 必为 0，故 \(p(x)\) 以 \((x-a)\) 为因子。
\end{proof}

命题 \ref{pro:9.9} 给出了小次数多项式不可约性的一个判别法：

\begin{proposition}\label{pro:9.10}
域 \(F\) 上一个二次或三次多项式可约当且仅当它在 \(F\) 中有根。
\end{proposition}

\begin{proof}
这由前一命题立即推出，因为二次或三次多项式可约当且仅当它至少有一个线性因子。
\end{proof}

下一个结果限制了整系数多项式根的可能性（为方便起见它是对 \(\mathbb{Z}[x]\) 叙述的，虽然它显然可以推广到 \(R[x]\)，其中 \(R\) 是任意唯一分解整环）。

\begin{proposition}\label{pro:9.11}
设 \(p(x) = a_nx^n + a_{n-1}x^{n-1} + \cdots + a_0\) 是 \(n\) 次整系数多项式。若 \(r/s \in \mathbb{Q}\) 是既约的（即 \(r\) 与 \(s\) 是互素整数）且 \(r/s\) 是 \(p(x)\) 的根，则 \(r\) 整除常数项而 \(s\) 整除 \(p(x)\) 的首项系数：\(r \mid a_0\)，\(s \mid a_n\). 特别地，若 \(p(x)\) 是整系数首一多项式，且对整除 \(p(x)\) 常数项的所有整数 \(d\) 都有 \(p(d) \neq 0\)，则 \(p(x)\) 在 \(\mathbb{Q}\) 中没有根。
\end{proposition}

\begin{proof}
由假设，\(p(r/s) = 0 = a_n(r/s)^n + a_{n-1}(r/s)^{n-1} + \cdots + a_0\). 两边乘以 \(s^n\)，得到
\[
a_nr^n + a_{n-1}r^{n-1}s + \cdots + a_1rs^{n-1} + a_0s^n = 0.
\]
于是 \(a_nr^n = s(-a_{n-1}r^{n-1} - \cdots - a_0s^{n-1})\)，故 \(s\) 整除 \(a_nr^n\). 由假设，\(s\) 与 \(r\) 互素，由此可得 \(s \mid a_n\). 类似地，解出方程中的 \(a_0s^n\) 可得 \(r \mid a_0\). 命题的最后一个断言由前面这些推出。
\end{proof}

\begin{example}
（1）多项式 \(x^3-3x-1\) 在 \(\mathbb{Z}[x]\) 中不可约。为证明这一点，由高斯引理和命题 \ref{pro:9.10}，只需证明它没有有理根。由命题 \ref{pro:9.11}，有理根的唯一候选是整除常数项 1 的整数，即 \(\pm 1\). 把 1 和 \(-1\) 代入多项式表明它们都不是根。

（2）对任意素数 \(p\)，多项式 \(x^2-p\) 和 \(x^3-p\) 在 \(\mathbb{Q}[x]\) 中不可约。这是因为它们的次数是 2 和 3，故只需证明它们没有有理根。由命题 \ref{pro:9.11}，根的唯一候选是 \(\pm 1\) 与 \(\pm p\)，而把它们代入多项式都不给出 0.

（3）多项式 \(x^2+1\) 在 \(\mathbb{Z}/2\mathbb{Z}[x]\) 中可约，因为 1 是它的根，且它分解为 \((x+1)^2\).

（4）多项式 \(x^2+x+1\) 在 \(\mathbb{Z}/2\mathbb{Z}[x]\) 中不可约，因为它在 \(\mathbb{Z}/2\mathbb{Z}\) 中没有根：\(0^2+0+1 = 1\) 且 \(1^2+1+1 = 1\).

（5）类似地，多项式 \(x^3+x+1\) 在 \(\mathbb{Z}/2\mathbb{Z}[x]\) 中不可约。
\end{example}

这个技巧限于低次多项式，因为它依赖于一次因子的存在。例如，四次多项式可能是两个不可约二次式之积，从而可约但没有线性因子。

一个相当一般的检验不可约性的技巧使用上面的命题 \ref{pro:9.2}，做法是把系数模某个理想约化。

\begin{proposition}\label{pro:9.12}
设 \(I\) 是整环 \(R\) 的真理想，\(p(x)\) 是 \(R[x]\) 中的非常数首一多项式。若 \(p(x)\) 在 \((R/I)[x]\) 中的像不能分解为 \((R/I)[x]\) 中两个次数更小的多项式之积，则 \(p(x)\) 在 \(R[x]\) 中不可约。
\end{proposition}

\begin{proof}
假设 \(p(x)\) 在 \((R/I)[x]\) 中不能分解，但 \(p(x)\) 在 \(R[x]\) 中可约。如上一节末尾所指出的，这意味着存在 \(R[x]\) 中首一的非常数多项式 \(a(x)\) 和 \(b(x)\) 使得 \(p(x) = a(x)b(x)\). 由命题 \ref{pro:9.2}，把系数模 \(I\) 约化就得到 \((R/I)[x]\) 中一个因子都非常数的分解，矛盾。
\end{proof}

这个命题表明，若能找到真理想 \(I\) 使得约化后的多项式不能分解，则该多项式本身不可约。遗憾的是，甚至在 \(\mathbb{Z}[x]\) 中也有这样的例子：多项式不可约，但它在模任何理想下的约化都可约（故它的不可约性无法用这个技巧检测出来）。例如，多项式 \(x^4+1\) 在 \(\mathbb{Z}[x]\) 中不可约，但模每个素数都可约（我们将在第 14 章验证这一点），而多项式 \(x^4-72x^2+4\) 在 \(\mathbb{Z}[x]\) 中不可约，但模每个整数都可约。

\begin{example}
（1）考虑 \(\mathbb{Z}[x]\) 中的多项式 \(p(x) = x^2+x+1\). 模 2 约化，由上面的例 4 我们看到 \(p(x)\) 在 \(\mathbb{Z}[x]\) 中不可约。类似地，\(x^3+x+1\) 在 \(\mathbb{Z}[x]\) 中不可约，因为它在 \(\mathbb{Z}/2\mathbb{Z}[x]\) 中不可约。

（2）多项式 \(x^2+1\) 在 \(\mathbb{Z}[x]\) 中不可约，因为它在 \(\mathbb{Z}/3\mathbb{Z}[x]\) 中不可约（在 \(\mathbb{Z}/3\mathbb{Z}\) 中没有根），但它模 2 可约。这表明命题 \ref{pro:9.12} 的逆不成立。

（3）模一个理想约化以判定不可约性的想法在多变量中也可以使用，但必须小心。例如，多项式 \(x^2+xy+1\) 在 \(\mathbb{Z}[x, y]\) 中不可约，因为模理想 \((y)\) 它化为 \(\mathbb{Z}[x]\) 中的 \(x^2+1\)，而后者不可约且次数相同。在这类论证中必须注意“坍缩”。例如，多项式 \(xy+x+y+1\)（它等于 \((x+1)(y+1)\)）可约，但它模 \((x)\) 和模 \((y)\) 看起来都不可约。原因是 \(\mathbb{Z}[x, y]\) 中的非单位多项式可以约化为商环中的单位。要顾及这一点，必须确定原环中哪些元素在商环中成为单位。例如，\(\mathbb{Z}[x, y]\) 中模 \((y)\) 为单位元素的多项式，正是那些常数项为 \(\pm 1\) 且所有非常数项都被 \(y\) 整除的 \(\mathbb{Z}[x, y]\) 中的多项式。\(x^2+xy+1\) 与它模 \((y)\) 的约化有相同的次数，这一事实就排除了出现一个模 \((y)\) 为单位、但在 \(\mathbb{Z}[x, y]\) 中不是单位的因子的可能性，从而给出这个多项式的不可约性。
\end{example}

把系数模一个理想约化以检验不可约性的一个特殊情形非常有用，它称为艾森斯坦判别法（虽然舍内曼（Schönemann）更早就证明了它，因此更恰当的名称是艾森斯坦--舍内曼判别法）：

\begin{proposition}\label{pro:9.13}
（\textbf{艾森斯坦判别法}）设 \(P\) 是整环 \(R\) 的素理想，\(f(x) = x^n + a_{n-1}x^{n-1} + \cdots + a_1x + a_0\) 是 \(R[x]\) 中的多项式（这里 \(n \geq 1\)）。假设 \(a_{n-1}, \ldots, a_1, a_0\) 全都是 \(P\) 的元素，且假设 \(a_0\) 不是 \(P^2\) 的元素。则 \(f(x)\) 在 \(R[x]\) 中不可约。
\end{proposition}

\begin{proof}
假设 \(f(x)\) 可约，比如在 \(R[x]\) 中 \(f(x) = a(x)b(x)\)，其中 \(a(x)\) 和 \(b(x)\) 是非常数多项式。把这个方程模 \(P\) 约化，并利用关于 \(f(x)\) 系数的假设，我们得到 \((R/P)[x]\) 中的方程 \(\bar{x}^n = \bar{a}(x)\bar{b}(x)\)，其中横线表示系数模 \(P\) 约化后的多项式。由于 \(P\) 是素理想，\(R/P\) 是整环，由此 \(a(x)\) 和 \(b(x)\) 的常数项都为 0，即 \(a(x)\) 和 \(b(x)\) 的常数项都是 \(P\) 的元素。但此时 \(f(x)\) 的常数项 \(a_0\) 作为这两个常数项之积将是 \(P^2\) 的元素，矛盾。
\end{proof}

艾森斯坦判别法最常应用于 \(\mathbb{Z}[x]\)，所以我们把这个情形显式写出：

\begin{corollary}\label{cor:9.14}
（\textbf{\(\mathbb{Z}[x]\) 上的艾森斯坦判别法}）设 \(p\) 是 \(\mathbb{Z}\) 中的素数，\(f(x) = x^n + a_{n-1}x^{n-1} + \cdots + a_1x + a_0 \in \mathbb{Z}[x]\)，\(n \geq 1\). 假设对每个 \(i \in \{0, 1, \ldots, n-1\}\) 都有 \(p \mid a_i\)，但 \(p^2\) 不整除 \(a_0\). 则 \(f(x)\) 在 \(\mathbb{Z}[x]\) 和 \(\mathbb{Q}[x]\) 中都不可约。
\end{corollary}

\begin{proof}
这只是命题 \ref{pro:9.13} 在 \(\mathbb{Z}\) 的素理想 \((p)\) 这一情形下的重述，再加上推论 \ref{cor:9.6}.
\end{proof}

\begin{example}
（1）多项式 \(x^4+10x+5\) 在 \(\mathbb{Z}[x]\) 中由对素数 5 应用艾森斯坦判别法可知不可约。

（2）若整数 \(a\) 被某个素数 \(p\) 整除但不被 \(p^2\) 整除，则 \(x^n-a\) 由艾森斯坦判别法在 \(\mathbb{Z}[x]\) 中不可约。特别地，对所有正整数 \(n\)，\(x^n-p\) 都不可约，于是对 \(n \geq 2\)，\(p\) 的 \(n\) 次根不是有理数（即这个多项式在 \(\mathbb{Q}\) 中没有根）。

（3）考虑前面提到的多项式 \(f(x) = x^4+1\). 艾森斯坦判别法不能直接用于 \(f(x)\). 多项式 \(g(x) = f(x+1)\) 是 \((x+1)^4+1\)，即 \(x^4+4x^3+6x^2+4x+2\)，对素数 2 应用艾森斯坦判别法表明这个多项式不可约。由此 \(f(x)\) 也必不可约，因为 \(f(x)\) 的任何分解都会给出 \(g(x)\) 的分解（只需把 \(x\) 换成 \(x+1\)）。这个例子表明艾森斯坦判别法有时可用来验证一个它不能直接应用的多项式的不可约性。

（4）作为另一个例子，设 \(p\) 是素数并考虑多项式
\[
\Phi_p(x) = \frac{x^p-1}{x-1} = x^{p-1} + x^{p-2} + \cdots + x + 1,
\]
它是分圆多项式的一个例子，我们将在第 IV 部分更充分地讨论。艾森斯坦判别法也不能直接应用，但它对素数 \(p\) 适用于多项式
\[
\Phi_p(x+1) = \frac{(x+1)^p-1}{x} = x^{p-1} + px^{p-2} + \cdots + \frac{p(p-1)}{2}x + p \in \mathbb{Z}[x],
\]
因为由二项式定理，除第一项外所有系数都被 \(p\) 整除。和前面一样，这说明 \(\Phi_p(x)\) 在 \(\mathbb{Z}[x]\) 中不可约。

（5）作为使用命题 \ref{pro:9.13} 中更一般形式的艾森斯坦判别法的例子，我们模仿上面的例 2. 设 \(R = \mathbb{Q}[x]\)，\(n\) 是任意正整数。考虑环 \(R[X]\) 中的多项式 \(X^n-x\). 理想 \((x)\) 是系数环 \(R\) 中的素理想，因为 \(R/(x) = \mathbb{Q}[x]/(x)\) 是整环 \(\mathbb{Q}\). 对 \(R\) 的理想 \((x)\) 应用艾森斯坦判别法直接表明 \(X^n-x\) 在 \(R[X]\) 中不可约。注意这个构造把 \(\mathbb{Q}\) 换成任何域、甚至任何整环都有效。
\end{example}

现在已有对某些域上的多项式进行分解的有效算法。对整系数多项式，这些算法已在若干计算机软件包中实现。对 \(\mathbb{F}_p\) 上多项式分解的一个有效算法称为 Berlekamp 算法，它将在第 14.3 节末的习题中详细描述。

\subsection*{习题}

\begin{enumerate}
\item 判断下列多项式在所指出的环中是否不可约。对可约的，确定它们的不可约分解。记号 \(\mathbb{F}_p\) 表示有限域 \(\mathbb{Z}/p\mathbb{Z}\)（\(p\) 是素数）。

（a）\(x^2+x+1\) 在 \(\mathbb{F}_2[x]\) 中。

（b）\(x^3+x+1\) 在 \(\mathbb{F}_3[x]\) 中。

（c）\(x^4+1\) 在 \(\mathbb{F}_5[x]\) 中。

（d）\(x^4+10x^2+1\) 在 \(\mathbb{Z}[x]\) 中。

\item 证明下列多项式在 \(\mathbb{Z}[x]\) 中不可约：

（a）\(x^4-4x^3+6\)

（b）\(x^6+30x^5-15x^3+6x-120\)

（c）\(x^4+4x^3+6x^2+2x+1\)〔把 \(x\) 换成 \(x-1\).〕

（d）\(\dfrac{(x+2)^p-2^p}{x}\)，其中 \(p\) 是奇素数。

\item 证明多项式 \((x-1)(x-2)\cdots(x-n)-1\) 对所有 \(n \geq 1\) 在 \(\mathbb{Z}\) 上不可约。〔若该多项式有分解，考虑各因子在 \(x = 1, 2, \ldots, n\) 处的值。〕

\item 证明多项式 \((x-1)(x-2)\cdots(x-n)+1\) 对所有 \(n \geq 1\)、\(n \neq 4\) 在 \(\mathbb{Z}\) 上不可约。

\item 求 \(\mathbb{F}_2[x]\) 中所有次数 \(\leq 3\) 的首一不可约多项式，并求出 \(\mathbb{F}_3[x]\) 中同样的多项式。

\item 构造下列各阶的域：（a）9，（b）49，（c）8，（d）81（你可以把它们表示为某个 \(F\) 和 \(f\) 的 \(F[x]/(f(x))\)）。〔利用第 9.2 节习题 2 和 3。〕

\item 证明 \(\mathbb{R}[x]/(x^2+1)\) 是域，且同构于复数。

\item 证明 \(K_1 = \mathbb{F}_{11}[x]/(x^2+1)\) 和 \(K_2 = \mathbb{F}_{11}[y]/(y^2+2y+2)\) 都是有 121 个元素的域。证明把 \(K_1\) 的元素 \(p(i)\) 映到 \(K_2\) 的元素 \(p(y+1)\) 的映射（其中 \(p\) 是任意系数在 \(\mathbb{F}_{11}\) 中的多项式）有定义，且给出从 \(K_1\) 到 \(K_2\) 的环（因而域）同构。

\item 证明多项式 \(x^2-\sqrt{2}\) 在 \(\mathbb{Z}[\sqrt{2}]\) 上不可约（你可以使用 \(\mathbb{Z}[\sqrt{2}]\) 是 U.F.D. 这一事实——参见第 8.1 节习题 9）。

\item 证明多项式 \(p(x) = x^4-4x^2+8x+2\) 在二次域 \(F = \mathbb{Q}(\sqrt{-2}) = \{a+b\sqrt{-2} \mid a, b \in \mathbb{Q}\}\) 上不可约。〔首先对唯一分解整环 \(\mathbb{Z}[\sqrt{-2}]\) 使用命题 \ref{pro:9.11} 的方法（参见第 8.1 节习题 8），证明若 \(\alpha \in \mathbb{Z}[\sqrt{-2}]\) 是 \(p(x)\) 的根，则 \(\alpha\) 是 2 在 \(\mathbb{Z}[\sqrt{-2}]\) 中的因子。由此断定 \(\alpha\) 必为 \(\pm 1\)、\(\pm\sqrt{-2}\) 或 \(\pm 2\)，从而证明 \(p(x)\) 在 \(F\) 上没有线性因子。类似地证明 \(p(x)\) 不是两个系数在 \(F\) 中的二次式之积。〕

\item 证明 \(x^2+y^2-1\) 在 \(\mathbb{Q}[x, y]\) 中不可约。

\item 证明 \(x^{n-1}+x^{n-2}+\cdots+x+1\) 在 \(\mathbb{Z}\) 上不可约当且仅当 \(n\) 是素数。

\item 证明 \(x^3+nx+2\) 对所有整数 \(n \neq 1, -3, -5\) 在 \(\mathbb{Z}\) 上不可约。

\item 把 \(x^8-1\) 和 \(x^6-1\) 这两个多项式分别在下列各环上分解为不可约因子：（a）\(\mathbb{Z}\)，（b）\(\mathbb{Z}/2\mathbb{Z}\)，（c）\(\mathbb{Z}/3\mathbb{Z}\).

\item 证明若 \(F\) 是域，则多项式 \(x^n-x\)（它的系数在形式幂级数环 \(F[[x]]\) 中，参见第 7.2 节习题 3）在 \(F[[x]]\) 上不可约。〔回忆 \(F[[x]]\) 是欧几里得整环——参见第 7.2 节习题 5 和第 8.1 节例 4。〕

\item 设 \(F\) 是域，\(f(x)\) 是 \(F[x]\) 中次数为 \(n\) 的多项式。多项式 \(g(x) = x^nf(1/x)\) 称为 \(f(x)\) 的\textbf{反多项式}。

（a）用 \(f\) 的系数描述 \(g\) 的系数。

（b）证明 \(f\) 不可约当且仅当 \(g\) 不可约。

\item 证明艾森斯坦判别法的下述变形：设 \(P\) 是唯一分解整环 \(R\) 中的素理想，\(f(x) = a_nx^n + a_{n-1}x^{n-1} + \cdots + a_1x + a_0\) 是 \(R[x]\) 中的多项式，\(n \geq 2\). 假设 \(a_n \notin P\)，\(a_{n-1}, \ldots, a_0 \in P\) 且 \(a_0 \notin P^2\). 证明 \(f(x)\) 在 \(F[x]\) 中不可约，其中 \(F\) 是 \(R\) 的商域。

\item 证明 \(6x^5+14x^3-21x+35\) 和 \(18x^5-30x^2+120x+360\) 在 \(\mathbb{Q}[x]\) 中不可约。

\item 设 \(F\) 是域，\(f(x) = a_nx^n + a_{n-1}x^{n-1} + \cdots + a_0 \in F[x]\). \(f(x)\) 的\textbf{导数} \(D_x(f(x))\) 定义为
\[
D_x(f(x)) = na_nx^{n-1} + (n-1)a_{n-1}x^{n-2} + \cdots + a_1,
\]
其中照常 \(na = a + a + \cdots + a\)（\(n\) 次）。注意 \(D_x(f(x))\) 仍是系数在 \(F\) 中的多项式。

称多项式 \(f(x)\) 有一个\textbf{重根}，若存在包含 \(F\) 的域 \(E\) 与某个 \(\alpha \in E\) 使得 \((x-\alpha)^2\) 在 \(E[x]\) 中整除 \(f(x)\). 例如，多项式 \(f(x) = (x-1)^2(x-2) \in \mathbb{Q}[x]\) 以 1 为重根，而多项式 \(f(x) = x^4+2x^2+1 = (x^2+1)^2 \in \mathbb{R}[x]\) 以 \(\pm i \in \mathbb{C}\) 为重根。我们将在第 13.5 节证明：非常数多项式 \(f(x)\) 有重根当且仅当 \(f(x)\) 与它的导数不互素（这可以用 \(F[x]\) 中的欧几里得算法检测）。用这个判别法判断下列多项式是否有重根：

（a）\(x^3-3x-2 \in \mathbb{Q}[x]\)

（b）\(x^3+3x+2 \in \mathbb{Q}[x]\)

（c）\(x^6-4x^4+6x^3+4x^2-12x+9 \in \mathbb{Q}[x]\)

（d）证明对任意素数 \(p\) 和任意 \(a \in \mathbb{F}_p\)，多项式 \(x^p-a\) 有重根。

\item 证明多项式 \(f(x) = x\) 在 \(\mathbb{Z}/6\mathbb{Z}[x]\) 中分解为 \((3x+4)(4x+3)\)，因而它不是不可约多项式。

（a）证明 \(f(x)\) 模 \(\mathbb{Z}/6\mathbb{Z}\) 的两个非平凡理想 \((2)\) 和 \((3)\) 的约化都是不可约多项式，这表明命题 \ref{pro:9.12} 中 \(R\) 是整环这一条件是必要的。

（b）证明在 \(\mathbb{Z}/6\mathbb{Z}[x]\) 的任意分解 \(f(x) = g(x)h(x)\) 中，\(g(x)\) 模 \((2)\) 的约化是 1 或 \(x\)，而 \(h(x)\) 模 \((2)\) 的约化相应地为 \(x\) 或 1；对模 \((3)\) 的约化也有同样的结论。确定 \(f(x)\) 在 \(\mathbb{Z}/6\mathbb{Z}[x]\) 中的全部分解。〔利用中国剩余定理。〕

（c）证明理想 \((3, x)\) 是 \(\mathbb{Z}/6\mathbb{Z}[x]\) 中的主理想。

（d）证明在环 \(\mathbb{Z}/30\mathbb{Z}[x]\) 上，多项式 \(f(x) = x\) 有分解 \(f(x) = (10x+21)(15x+16)(6x+25)\). 证明这些因子的任意两个之积仍与 \(f(x)\) 有相同的次数。证明 \(f(x)\) 模 \(\mathbb{Z}/30\mathbb{Z}\) 中任意素数的约化都是不可约多项式。确定 \(f(x)\) 在 \(\mathbb{Z}/30\mathbb{Z}[x]\) 中的全部分解。〔考虑模 \((2)\)、\((3)\) 和 \((5)\) 的约化，并利用中国剩余定理。〕

（e）把（d）推广到 \(\mathbb{Z}/n\mathbb{Z}[x]\)，其中 \(n\) 是 \(k\) 个互异素数之积。
\end{enumerate}

\section{域上的多项式环（二）}

设 \(F\) 是域。这里我们证明关于一个变量多项式环 \(F[x]\) 的一些补充结果。第一个是先前所得到结果的重述。

\begin{proposition}\label{pro:9.15}
\(F[x]\) 中的极大理想正是由不可约多项式 \(f(x)\) 生成的理想 \((f(x))\). 特别地，\(F[x]/(f(x))\) 是域当且仅当 \(f(x)\) 不可约。
\end{proposition}

\begin{proof}
这由第 8.2 节命题 7 应用于主理想整环 \(F[x]\) 得到。
\end{proof}

\begin{proposition}\label{pro:9.16}
设 \(g(x)\) 是 \(F[x]\) 中的非常数元素，并设
\[
g(x) = f_1(x)^{n_1}f_2(x)^{n_2}\cdots f_k(x)^{n_k}
\]
是它分解为不可约元素的分解，其中 \(f_i(x)\) 互不相同。则我们有下列环同构：
\[
F[x]/(g(x)) \cong F[x]/(f_1(x)^{n_1}) \times F[x]/(f_2(x)^{n_2}) \times \cdots \times F[x]/(f_k(x)^{n_k}).
\]
\end{proposition}

\begin{proof}
这由中国剩余定理（定理 7.17）推出，因为当 \(f_i(x)\) 与 \(f_j(x)\) 互不相同时，理想 \((f_i(x)^{n_i})\) 与 \((f_j(x)^{n_j})\) 互余（它们在欧几里得整环 \(F[x]\) 中互素，从而由它们生成的理想是 \(F[x]\)）。
\end{proof}

下一个结果涉及域 \(F\) 上多项式根的个数。由命题 \ref{pro:9.9}，一个根 \(\alpha\) 对应于 \(f(x)\) 的一个线性因子 \((x-\alpha)\). 若 \(f(x)\) 被 \((x-\alpha)^m\) 整除但不被 \((x-\alpha)^{m+1}\) 整除，则称 \(\alpha\) 是 \(m\) 重根。

\begin{proposition}\label{pro:9.17}
若多项式 \(f(x)\) 在 \(F\) 中有根 \(\alpha_1, \alpha_2, \ldots, \alpha_k\)（不必互异），则 \(f(x)\) 以 \((x-\alpha_1)\cdots(x-\alpha_k)\) 为因子。特别地，域 \(F\) 上一个变量的 \(n\) 次多项式在 \(F\) 中至多有 \(n\) 个根，即使按重数计也是如此。
\end{proposition}

\begin{proof}
第一个断言由命题 \ref{pro:9.9} 简单归纳即得。由于线性因子不可约，第二个断言由 \(F[x]\) 是唯一分解整环推出。
\end{proof}

这最后一个结果有下面这个有意思的推论。

\begin{proposition}\label{pro:9.18}
域的乘法群的有限子群是循环群。特别地，若 \(F\) 是有限域，则 \(F\) 的非零元素乘法群 \(F^\times\) 是循环群。
\end{proposition}

\begin{proof}
我们用有限生成阿贝尔群基本定理（第 5.2 节定理 3）给出一个证明。习题中给出一个更数论化的证明，或者也可以用第 6.1 节命题 5 代替基本定理。由基本定理，这个有限子群可以写成循环群的直积
\[
\mathbb{Z}/n_1\mathbb{Z} \times \mathbb{Z}/n_2\mathbb{Z} \times \cdots \times \mathbb{Z}/n_k\mathbb{Z},
\]
其中 \(n_k \mid n_{k-1} \mid \cdots \mid n_2 \mid n_1\). 一般地，若 \(G\) 是循环群且 \(d \mid |G|\)，则 \(G\) 中恰有 \(d\) 个阶整除 \(d\) 的元素。由于 \(n_k\) 整除这个直积中每个循环群的阶，可知每个直因子都含有 \(n_k\) 个阶整除 \(n_k\) 的元素。若 \(k > 1\)，这样的元素总数就超过 \(n_k\). 但那样域 \(F\) 中多项式 \(x^{n_k}-1\) 就有多于 \(n_k\) 个根，与命题 \ref{pro:9.17} 矛盾。故 \(k = 1\)，这个群是循环群。
\end{proof}

\begin{corollary}\label{cor:9.19}
设 \(p\) 是素数。非零模 \(p\) 剩余类的乘法群 \((\mathbb{Z}/p\mathbb{Z})^\times\) 是循环群。
\end{corollary}

\begin{proof}
这是有限域 \(\mathbb{Z}/p\mathbb{Z}\) 的乘法群。
\end{proof}

\begin{corollary}\label{cor:9.20}
设 \(n \geq 2\) 是整数，在 \(\mathbb{Z}\) 中的分解为 \(n = p_1^{\alpha_1}p_2^{\alpha_2}\cdots p_r^{\alpha_r}\)，其中 \(p_1, \ldots, p_r\) 是互不相同的素数。我们有下列（乘法）群的同构：

（1）\((\mathbb{Z}/n\mathbb{Z})^\times \cong (\mathbb{Z}/p_1^{\alpha_1}\mathbb{Z})^\times \times (\mathbb{Z}/p_2^{\alpha_2}\mathbb{Z})^\times \times \cdots \times (\mathbb{Z}/p_r^{\alpha_r}\mathbb{Z})^\times\)

（2）\((\mathbb{Z}/2^\alpha\mathbb{Z})^\times\) 是一个 2 阶循环群与一个 \(2^{\alpha-2}\) 阶循环群的直积，对所有 \(\alpha \geq 2\)；

（3）\((\mathbb{Z}/p^\alpha\mathbb{Z})^\times\) 是 \(p^{\alpha-1}(p-1)\) 阶循环群，对所有奇素数 \(p\).
\end{corollary}

\noindent\textbf{注：}这些同构描述了 \(n\) 阶循环群 \(Z_n\) 的自同构群的群论结构，因为 \(\operatorname{Aut}(Z_n) \cong (\mathbb{Z}/n\mathbb{Z})^\times\)（参见第 4.4 节命题 16）。特别地，对素数 \(p\)，\(p\) 阶循环群的自同构群是 \(p-1\) 阶循环群。

\begin{proof}
这主要是汇集先前的结果。第（1）个同构由中国剩余定理推出（参见第 7.6 节推论 18）。（2）中的同构直接由第 2.3 节习题 22 和 23 推出。

对奇素数 \(p\)，\((\mathbb{Z}/p^\alpha\mathbb{Z})^\times\) 是 \(p^{\alpha-1}(p-1)\) 阶阿贝尔群。由第 2.3 节习题 21，这个群的西罗 \(p\)-子群是循环的。由
\[
\mathbb{Z}/p^\alpha\mathbb{Z} \longrightarrow \mathbb{Z}/p\mathbb{Z} \qquad (\text{模 } p \text{ 约化})
\]
定义的映射是环同态，它给出从 \((\mathbb{Z}/p^\alpha\mathbb{Z})^\times\) 到 \((\mathbb{Z}/p\mathbb{Z})^\times\) 的满群同态。后一个群是 \(p-1\) 阶循环群（推论 \ref{cor:9.19}）。这个映射的核的阶是 \(p^{\alpha-1}\)，故对所有素数 \(q \neq p\)，\((\mathbb{Z}/p^\alpha\mathbb{Z})^\times\) 的西罗 \(q\)-子群同构地映入循环群 \((\mathbb{Z}/p\mathbb{Z})^\times\). 于是 \((\mathbb{Z}/p^\alpha\mathbb{Z})^\times\) 的所有西罗子群都是循环的，故（3）成立，证明完成。
\end{proof}

\subsection*{习题}

\begin{enumerate}
\item 设 \(F\) 是域，\(f(x)\) 是 \(F[x]\) 中的非常数多项式。用 \(f(x)\) 的分解描述 \(F[x]/(f(x))\) 的幂零根（参见第 7.3 节习题 29）。

\item 对第 9.4 节习题 6 中构造的每个域，给出其（循环的）非零元素乘法群的一个生成元。

\item 设 \(p\) 是 \(\mathbb{Z}\) 中的奇素数，\(n\) 是正整数。证明 \(x^n-p\) 在 \(\mathbb{Z}[i]\) 上不可约。〔利用第 8 章的命题 18 与艾森斯坦判别法。〕

\item 证明 \(x^3+12x^2+18x+6\) 在 \(\mathbb{Z}[i]\) 上不可约。〔利用命题 8.18 与艾森斯坦判别法。〕

\item 设 \(\varphi\) 表示欧拉 \(\varphi\)-函数。证明恒等式 \(\sum_{d \mid n}\varphi(d) = n\)，其中求和遍历 \(n\) 的所有因子 \(d\). 〔首先注意当 \(n = p^m\) 是素数 \(p\) 的幂时这个恒等式成立，因为求和可逐项相消。写 \(n = p^mn'\)，其中 \(p\) 不整除 \(n'\). 通过把右端乘开并利用 \(a\) 与 \(b\) 互素时 \(\varphi(ab) = \varphi(a)\varphi(b)\) 的乘性，证明 \(\sum_{d \mid n}\varphi(d) = \sum_{d'' \mid p^m}\varphi(d'')\sum_{d' \mid n'}\varphi(d')\). 用归纳法完成证明。这个问题也可以这样处理：设 \(Z\) 是 \(n\) 阶循环群，并证明由于对每个整除 \(n\) 的 \(d\)，\(Z\) 都含有唯一的 \(d\) 阶子群，\(Z\) 中 \(d\) 阶元素的个数是 \(\varphi(d)\). 于是 \(|Z|\) 是 \(\varphi(d)\) 对 \(n\) 的所有因子 \(d\) 求和。〕

\item 设 \(G\) 是域 \(F\) 的非零元素乘法群 \(F^\times\) 的一个 \(n\) 阶有限子群。设 \(\varphi\) 表示欧拉 \(\varphi\)-函数，\(\psi(d)\) 表示 \(G\) 中 \(d\) 阶元素的个数。证明对 \(n\) 的每个因子 \(d\) 都有 \(\psi(d) = \varphi(d)\). 特别地断定 \(\psi(n) \geq 1\)，从而 \(G\) 是循环群。〔注意对任意整数 \(N \geq 1\)，多项式 \(x^N-1\) 在 \(F\) 中至多有 \(N\) 个根。由此断定对任意整数 \(N\) 有 \(\sum_{d \mid N}\psi(d) \leq N\). 由上一习题 \(\sum_{d \mid N}\varphi(d) = N\)，用归纳法证明对 \(n\) 的每个因子 \(d\) 有 \(\psi(d) \leq \varphi(d)\). 由于 \(\sum_{d \mid n}\psi(d) = n = \sum_{d \mid n}\varphi(d)\)，证明这蕴涵对 \(n\) 的每个因子 \(d\) 有 \(\psi(d) = \varphi(d)\).〕

\item 证明域的加法群与乘法群绝不同构。〔考虑三种情形：\(|F|\) 有限；\(F\) 中 \(-1 \neq 1\)；\(F\) 中 \(-1 = 1\).〕
\end{enumerate}

\section{域上多变量多项式与 Gröbner 基}

本节我们考虑多变量多项式，介绍一些基本的计算工具，并指出一些应用。本节的结果在第 10 章到第 14 章中并不需要。更多的应用将在第 15 章给出。

我们在第 2 节证明了域 \(F\) 上变量 \(x\) 的多项式环 \(F[x]\) 是欧几里得整环，推论 \ref{cor:9.8} 表明多项式环 \(F[x_1, \ldots, x_n]\) 是 U.F.D. 然而由第 8.2 节推论 8 可知，除非 \(n = 1\)，后面这个环不是 P.I.D. 下面第一个结果表明，这种多项式环中的理想虽然不必是主理想，却总是有限生成的。具有这个性质的一般环有一个专门的名称：

\begin{definition}
带有 1 的交换环 \(R\) 称为\textbf{诺特环}，若 \(R\) 的每个理想都是有限生成的。
\end{definition}

诺特环将在第 15 章和第 16 章更详细地研究。本节我们发展一些基本理论以及由此得到的算法，用以处理 \(F[x_1, \ldots, x_n]\) 中的（有限生成的）理想。

正如我们在第 1 节看到的，\(n\) 个变量的多项式环可以看作一个变量的、系数在 \(n-1\) 个变量的多项式环中的多项式环。沿着这种归纳方式——正如我们在定理 \ref{thm:9.7} 和推论 \ref{cor:9.8} 中所做的——我们可以由下面这个更一般的结果推出 \(F[x_1, x_2, \ldots, x_n]\) 是诺特环。

\begin{theorem}\label{thm:9.21}
（\textbf{希尔伯特基定理}）若 \(R\) 是诺特环，则多项式环 \(R[x]\) 也是诺特环。
\end{theorem}

\begin{proof}
设 \(I\) 是 \(R[x]\) 中的理想，\(L\) 是 \(I\) 中所有元素的首项系数组成的集合。我们先证明 \(L\) 是 \(R\) 的理想，如下。由于 \(I\) 含零多项式，\(0 \in L\). 设 \(f = ax^d + \cdots\) 和 \(g = bx^e + \cdots\) 是 \(I\) 中次数为 \(d, e\)、首项系数为 \(a, b \in R\) 的多项式。则对任意 \(r \in R\)，\(ra - b\) 要么是零，要么是多项式 \(rxf - x^dg\) 的首项系数。由于后一个多项式属于 \(I\)，我们有 \(ra - b \in L\)，这表明 \(L\) 是 \(R\) 的理想。由于假设 \(R\) 是诺特环，\(L\) 中的理想 \(L\) 是有限生成的，比如说由 \(a_1, a_2, \ldots, a_n \in R\) 生成。对每个 \(i = 1, \ldots, n\)，设 \(f_i\) 是 \(I\) 中首项系数为 \(a_i\) 的元素。设 \(e_i\) 表示 \(f_i\) 的次数，\(N\) 是 \(e_1, e_2, \ldots, e_n\) 的最大值。

对每个 \(d \in \{0, 1, \ldots, N-1\}\)，设 \(L_d\) 是 \(I\) 中次数为 \(d\) 的多项式的所有首项系数并上 0 组成的集合。与 \(L\) 类似的论证表明每个 \(L_d\) 也是 \(R\) 的理想，由于 \(R\) 是诺特环，它也是有限生成的。对每个非零理想 \(L_d\)，取 \(b_{d,1}, b_{d,2}, \ldots, b_{d,n_d} \in R\) 作为 \(L_d\) 的一组生成元，并设 \(f_{d,i}\) 是 \(I\) 中次数为 \(d\)、首项系数为 \(b_{d,i}\) 的多项式。

我们证明多项式 \(f_1, \ldots, f_n\) 连同所有非零理想 \(L_d\) 的所有多项式 \(f_{d,i}\) 构成 \(I\) 的一组生成元，即
\[
I = \left(\{f_1, \ldots, f_n\} \cup \{f_{d,i} \mid 0 \leq d < N,\ 1 \leq i \leq n_d\}\right).
\]
按构造，上面右端的理想 \(I'\) 包含于 \(I\)，因为所有生成元都是在 \(I\) 中选取的。若 \(I' \neq I\)，则存在 \(I\) 中不属于 \(I'\) 的次数最小的非零多项式 \(f\). 设 \(d = \deg f\)，\(a\) 是 \(f\) 的首项系数。

先假设 \(d \geq N\). 由于 \(a \in L\)，我们可以把 \(a\) 写成 \(L\) 的生成元的 \(R\)-线性组合：\(a = r_1a_1 + \cdots + r_na_n\). 则 \(g = r_1x^{d-e_1}f_1 + \cdots + r_nx^{d-e_n}f_n\) 是 \(I'\) 中与 \(f\) 有相同次数 \(d\) 和相同首项系数 \(a\) 的元素。于是 \(f - g \in I\) 是 \(I\) 中次数小于 \(f\) 的多项式。由 \(f\) 的最小性，必有 \(f - g = 0\)，故 \(f = g \in I'\)，矛盾。

再假设 \(d < N\). 此时对某个 \(d < N\) 有 \(a \in L_d\)，于是可以对某些 \(r_i \in R\) 写 \(a = r_1b_{d,1} + \cdots + r_{n_d}b_{d,n_d}\). 则 \(g = r_1f_{d,1} + \cdots + r_{n_d}f_{d,n_d}\) 是 \(I'\) 中与 \(f\) 有相同次数 \(d\) 和相同首项系数 \(a\) 的多项式，和前面一样得到矛盾。

由此 \(I = I'\) 是有限生成的，而由于 \(I\) 是任意的，这就完成了 \(R[x]\) 是诺特环的证明。
\end{proof}

由于域显然是诺特环，希尔伯特基定理与归纳立即给出：

\begin{corollary}\label{cor:9.22}
系数取自域 \(F\) 的多项式环 \(F[x_1, x_2, \ldots, x_n]\) 中的每个理想都是有限生成的。
\end{corollary}

若 \(I\) 是 \(F[x_1, \ldots, x_n]\) 中由多项式集合 \(S\)（可以无限）生成的理想，推论 \ref{cor:9.22} 表明 \(I\) 是有限生成的，事实上 \(I\) 由 \(S\) 中有限多个多项式生成（参见习题 1）。

正如希尔伯特基定理的证明所显示的，\(R[x]\) 中理想 \(I\) 的各多项式首项系数组成的集合是 \(R\) 中一个极有用的理想，可用来理解 \(I\). 这提示我们稍微内蕴地一般地研究 \(F[x_1, x_2, \ldots, x_n]\) 中的“首项”。为此我们需要指定单项式上的一个全序，因为若没有某种次序，一般无法判断一个多项式的“首项”是哪一个。在推论 \ref{cor:9.22} 的归纳证明中我们已经隐含地选择了这样一种次序——我们先把多项式 \(f\) 看作 \(x_1\) 的多项式，其系数在 \(R = F[x_2, \ldots, x_n]\) 中；然后把它的“首项系数”看作 \(x_2\) 的多项式，其系数在 \(F[x_3, \ldots, x_n]\) 中；依此类推。这是多项式环 \(F[x_1, \ldots, x_n]\) 上一个\textbf{字典单项序}的例子，它的定义方法是：先声明变量的次序，例如
\[
x_1 > x_2 > \cdots > x_n,
\]
然后规定指数为 \((a_1, a_2, \ldots, a_n)\) 的单项式项 \(Ax_1^{a_1}x_2^{a_2}\cdots x_n^{a_n}\) 的次序高于指数为 \((b_1, b_2, \ldots, b_n)\) 的单项式项 \(Bx_1^{b_1}x_2^{b_2}\cdots x_n^{b_n}\)，若这两个 \(n\) 元组首次出现不同的分量处有 \(a_i > b_i\). 这与字典中使用的次序类似（因此得名）：在字典中字母“a”在“b”之前，“b”又在“c”之前，等等，于是“aardvark”在“abacus”之前（尽管按字典序“单词”\(a_2a_3a_1\) 在 \(a_2a_1a_3\) 之前）。注意这个次序只在乘以单位（\(F^\times\) 的元素）的意义下定义，并且把两个单项式乘以同一个非零单项式不改变它们的次序。这可以一般地形式化如下。

\begin{definition}
\textbf{单项序}是单项式集合上的良序 \(\geq\)，它满足：对单项式 \(m, m_1, m_2\)，只要 \(m_1 \geq m_2\) 就有 \(mm_1 \geq mm_2\).

等价地，单项序也可以通过定义单项式多重次数 \(n\) 元组 \((a_1, \ldots, a_n) = \alpha\) 上的一个良序来给出，其中 \(m = Ax_1^{a_1}\cdots x_n^{a_n}\)，这个良序满足：若 \(\alpha \geq \beta\) 则 \(\alpha + \gamma \geq \beta + \gamma\).
\end{definition}

容易证明对任何单项序和任何单项式 \(m\) 都有 \(m \geq 1\)（参见习题 2）。利用希尔伯特基定理不难证明：单项式上任何全序，只要对每个单项式 \(m\) 满足 \(m \geq 1\)，并且只要 \(m_1 \geq m_2\) 就有 \(mm_1 \geq mm_2\)，就必定是良序（因而是单项序）——这组等价的单项序公理可能更容易验证。为简单起见，我们把例子限制在特别简单直观的字典序上，但重要的是要注意：在实践中使用其他单项序有重要的计算优势。习题中引入了一些其他常用的单项序。

如上所述，一旦有了单项序，我们就可以定义多项式的首项：

\begin{definition}
在多项式环 \(F[x_1, x_2, \ldots, x_n]\) 上固定一个单项序。

（1）\(F[x_1, x_2, \ldots, x_n]\) 中非零多项式 \(f\) 的\textbf{首项}记作 \(\operatorname{LT}(f)\)，是 \(f\) 中次序最大的单项式项，并规定 0 的首项为 0. 定义 \(f\) 的\textbf{多重次数}记作 \(\alpha(f)\)，是 \(f\) 的首项的多重次数。

（2）若 \(I\) 是 \(F[x_1, x_2, \ldots, x_n]\) 中的理想，则\textbf{首项理想}记作 \(\operatorname{LT}(I)\)，是由这个理想中所有元素的首项生成的理想，即
\[
\operatorname{LT}(I) = (\operatorname{LT}(f) \mid f \in I).
\]
\end{definition}

多项式的首项与多重次数显然依赖于次序的选择。例如若 \(x > y\)，则 \(\operatorname{LT}(2xy+y^3) = 2xy\)，多重次数为 \((1,1)\)；但若 \(y > x\)，则 \(\operatorname{LT}(2xy+y^3) = y^3\)，多重次数为 \((0,3)\). 特别地，多项式的首项不必是总次数最大的项。类似地，理想 \(I\) 的首项理想 \(\operatorname{LT}(I)\) 一般依赖于所用的次序。还要注意当 \(f\) 与 \(g\) 非零时，多项式的多重次数满足 \(\alpha(fg) = \alpha(f)+\alpha(g)\)，且此时 \(\operatorname{LT}(fg) = \operatorname{LT}(f)\operatorname{LT}(g)\)（参见习题 2）。

由定义，理想 \(\operatorname{LT}(I)\) 由单项式生成。这样的理想称为\textbf{单项式理想}，通常比一般的理想容易处理得多。例如，一个多项式属于单项式理想当且仅当它的每个单项式项都是该理想某个生成元的倍数（参见习题 10）。

在希尔伯特基定理的证明中，重要的是要有理想 \(I\) 的所有首项。若 \(I = (f_1, \ldots, f_m)\)，则按定义 \(\operatorname{LT}(I)\) 包含 \(I\) 的生成元的首项 \(\operatorname{LT}(f_1), \ldots, \operatorname{LT}(f_m)\). 由于 \(\operatorname{LT}(I)\) 是理想，它包含这些首项生成的理想：
\[
(\operatorname{LT}(f_1), \ldots, \operatorname{LT}(f_m)) \subseteq \operatorname{LT}(I).
\]
下面第一个例子表明，理想的首项理想 \(\operatorname{LT}(I)\) 一般可以严格大于仅由 \(I\) 的某组生成元的首项生成的理想。

\begin{example}
（1）在 \(F[x, y]\) 上选择字典序 \(x > y\). 多项式 \(f_1 = x^3y-xy^2+1\) 和 \(f_2 = x^2y^2-y^3-1\) 的首项是 \(\operatorname{LT}(f_1) = x^3y\)（故 \(f_1\) 的多重次数 \(\alpha(f_1) = (3,1)\)）和 \(\operatorname{LT}(f_2) = x^2y^2\)（故 \(\alpha(f_2) = (2,2)\)）。若 \(I = (f_1, f_2)\) 是由 \(f_1\) 和 \(f_2\) 生成的理想，则首项理想 \(\operatorname{LT}(I)\) 包含 \(\operatorname{LT}(f_1) = x^3y\) 和 \(\operatorname{LT}(f_2) = x^2y^2\)，故 \((x^3y, x^2y^2) \subseteq \operatorname{LT}(I)\). 由于
\[
yf_1 - xf_2 = y(x^3y-xy^2+1) - x(x^2y^2-y^3-1) = x+y,
\]
我们看到 \(g = x+y\) 是 \(I\) 的元素，故理想 \(\operatorname{LT}(I)\) 还包含首项 \(\operatorname{LT}(g) = x\). 这表明 \(\operatorname{LT}(I)\) 严格大于 \((\operatorname{LT}(f_1), \operatorname{LT}(f_2))\)，因为 \((\operatorname{LT}(f_1), \operatorname{LT}(f_2)) = (x^3y, x^2y^2)\) 中每个元素的总次数至少为 4. 我们将在后面看到，此时 \(\operatorname{LT}(I) = (x, y^4)\).

（2）关于字典序 \(y > x\)，上一例中 \(f_1\) 和 \(f_2\) 的首项是 \(\operatorname{LT}(f_1) = -xy^2\)（可以写成 \(-y^2x\) 以强调所选的次序）和 \(\operatorname{LT}(f_2) = -y^3\). 我们将在后面看到，在这个次序下 \(\operatorname{LT}(I) = (x^4, y)\)，这与上一例中使用次序 \(x > y\) 得到的首项理想不同，并且同样严格大于 \((\operatorname{LT}(f_1), \operatorname{LT}(f_2))\).

（3）在 \(F[x, y]\) 上选择任意次序，并设 \(f = f(x, y)\) 是任意非零多项式。此时主理想 \(I = (f)\) 中每个元素的首项都是 \(f\) 的首项的倍数，故此时 \(\operatorname{LT}(I) = (\operatorname{LT}(f))\).
\end{example}

在一个变量的情形，首项被用于除法算法中，把一个多项式 \(g\) 模另一个多项式 \(f\) 约化得到唯一的余式 \(r\)，而这个余式为 0 当且仅当 \(g\) 属于理想 \((f)\). 由于当 \(n \geq 2\) 时 \(F[x_1, x_2, \ldots, x_n]\) 不是欧几里得整环（因为它不是 P.I.D.），多变量多项式的情形更为复杂。在上面第一个例子中，\(f_1\) 和 \(f_2\) 都不整除 \(g\)（例如由次数考虑），故若试图先把 \(g\) 除以 \(f_1\) 或 \(f_2\) 之一、再除以另一个来把 \(g\) 模理想 \(I\) 约化，将会得到 \(g\) 本身的（非零）“余式”。特别地，这似乎表明 \(g = yf_1-xf_2\) 不是理想 \(I\) 的元素，尽管它确实是。次数为 1 的多项式 \(g\) 能成为两个次数为 4 的多项式 \(f_1\) 和 \(f_2\) 的线性组合，原因是在 \(yf_1\) 与 \(xf_2\) 之差中 \(y\) 的首项相消了，而这反映在 \(\operatorname{LT}(f_1)\) 和 \(\operatorname{LT}(f_2)\) 不足以生成 \(\operatorname{LT}(I)\) 这一事实中。\(F[x_1, \ldots, x_n]\) 中理想 \(I\) 的一组生成元，若其首项生成 \(I\) 中所有元素的首项，就有一个专门的名称。

\begin{definition}
多项式环 \(F[x_1, \ldots, x_n]\) 中理想 \(I\) 的 \textbf{Gröbner 基}是 \(I\) 的一个有限生成元集 \(\{g_1, \ldots, g_m\}\)，它的首项生成 \(I\) 中所有首项构成的理想，即
\[
(\operatorname{LT}(g_1), \ldots, \operatorname{LT}(g_m)) = \operatorname{LT}(I).
\]
\end{definition}

\noindent\textbf{注：}注意 Gröbner“基”实际上是 \(I\) 的一组生成元（它依赖于次序的选取），即 \(I\) 中每个元素都是这些生成元的线性组合，而不是向量空间意义下的基（那里线性组合是唯一的，参见第 10.3 节和第 11.1 节）。虽然可能有点容易引起误解，但“Gröbner 基”这个术语已被广泛采用，若引入别的名称反而有害。

Gröbner 基最重要的性质之一（将在后面的定理 \ref{thm:9.23} 中证明）是：每个多项式 \(g\) 都可以唯一地写成 \(I\) 中一个元素与一个由一般多项式除法得到的余式 \(r\) 之和。特别地，我们将看到 \(g\) 是 \(I\) 的元素当且仅当这个余式 \(r\) 为 0. 虽然一般来说也有类似的分解，但我们将看到若不使用 Gröbner 基，唯一性就丧失了（并且我们无法通过检验余式是否为 0 来判断是否属于 \(I\)），因为有些首项没有被生成元的首项所覆盖。

我们首先用 \(F[x_1, \ldots, x_n]\) 上由单项序定义的多项式的首项，把一个变量的除法算法推广为多变量中的非标准多项式除法。回忆对一个变量的多项式，通常的除法算法通过依次检验被除式 \(f(x)\) 的首项是否被 \(g(x)\) 的首项整除，来确定方程 \(f(x) = q(x)g(x)+r(x)\) 中的商 \(q(x)\) 与余式 \(r(x)\)：若 \(\operatorname{LT}(f) = a(x)\operatorname{LT}(g)\)，则把单项式项 \(a(x)\) 加到商上，并用被除式 \(f(x)-a(x)g(x)\)（由于首项相消——这是由 \(a(x)\) 的选取保证的——它的次数更小）代替 \(f(x)\) 重复这一过程。当除式 \(g(x)\) 的首项不再整除被除式的首项时过程终止，留下余式 \(r(x)\). 只需允许连续作除法，就可以把它推广为用有限多个多变量多项式作除法，得到一个余式和若干个商，如下。

\noindent\textbf{一般多项式除法}

在 \(F[x_1, \ldots, x_n]\) 上固定一个单项序，并设 \(g_1, \ldots, g_m\) 是 \(F[x_1, \ldots, x_n]\) 中一组非零多项式。若 \(f\) 是 \(F[x_1, \ldots, x_n]\) 中的任意多项式，开始时取商 \(q_1, \ldots, q_m\) 与余式 \(r\) 全为 0，并按 \(g_1, \ldots, g_m\) 的次序依次检验被除式 \(f\) 的首项是否被各除式的首项整除。然后

\begin{enumerate}
\item[(i)] 若 \(\operatorname{LT}(f)\) 被某个 \(\operatorname{LT}(g_i)\) 整除，比如 \(\operatorname{LT}(f) = a_i\operatorname{LT}(g_i)\)，则把 \(a_i\) 加到商 \(q_i\) 上，用被除式 \(f-a_ig_i\)（一个首项次序更低的多项式）代替 \(f\)，并重复整个过程。

\item[(ii)] 若被除式 \(f\) 的首项不被 \(\operatorname{LT}(g_1), \ldots, \operatorname{LT}(g_m)\) 中任何一个整除，则把 \(f\) 的首项加到余式 \(r\) 上，用被除式 \(f-\operatorname{LT}(f)\)（即去掉 \(f\) 的首项）代替 \(f\)，并重复整个过程。
\end{enumerate}

这个过程在有限步后终止（参见习题 3），此时被除式为 0，且得到一组商 \(q_1, \ldots, q_m\) 与一个余式 \(r\) 使得
\[
f = q_1g_1 + \cdots + q_mg_m + r.
\]
每个 \(q_ig_i\) 的多重次数都小于或等于 \(f\) 的多重次数，且余式 \(r\) 具有性质：\(r\) 中没有任何非零项被 \(\operatorname{LT}(g_1), \ldots, \operatorname{LT}(g_m)\) 中任何一个整除（因为在第（ii）步中只有具有这一性质的项才会被加到 \(r\) 上）。

\begin{example}
在 \(F[x, y]\) 上固定字典序 \(x > y\).

（1）设 \(f = x^3y^3+3x^2y^4\)，\(g = xy^4\). \(f\) 的首项是 \(x^3y^3\)，它不被 \(g\) 的（首项）整除，于是 \(x^3y^3\) 被加到余式 \(r\) 上（此时 \(r = x^3y^3\)），并用 \(f-\operatorname{LT}(f) = 3x^2y^4\) 代替 \(f\)，重新开始。由于 \(3x^2y^4\) 被 \(\operatorname{LT}(g) = xy^4\) 整除，商为 \(a = 3x\)，我们把 \(3x\) 加到商 \(q\) 上（此时 \(q = 3x\)），并用 \(3x^2y^4-a\operatorname{LT}(g) = 0\) 代替 \(3x^2y^4\)，此时过程终止。结果是商 \(q = 3x\)、余式 \(r = x^3y^3\)，且
\[
x^3y^3+3x^2y^4 = f = qg+r = (3x)(xy^4) + x^3y^3.
\]
注意若我们在第一步就因为 \(f\) 的首项不被 \(g\) 的首项整除而终止（这在一个变量的除法算法中是终止条件），我们就会剩下 \(f\) 本身作为“余式”，尽管 \(f\) 的“更多部分”是被 \(g\) 整除的。这就是除法过程中第（ii）步的原因（对一个变量的多项式这一步是不必要的）。

（2）设 \(f = x^2+x-y^2+y\)，并设 \(g_1 = xy+1\)、\(g_2 = x+y\). 第一次迭代中，\(f\) 的首项 \(x^2\) 不被 \(g_1\) 的首项整除，但被 \(g_2\) 的首项整除，故商 \(q_2\) 为 \(x\)，被除式 \(f\) 被替换为 \(f-xg_2 = -xy+x-y^2+y\). 第二次迭代中，\(-xy+x-y^2+y\) 的首项被 \(\operatorname{LT}(g_1)\) 整除，商为 \(-1\)，故 \(q_1 = -1\)，被除式被替换为 \((-xy+x-y^2+y)-(-1)g_1 = x-y^2+y+1\). 第三次迭代中，\(x-y^2+y+1\) 的首项不被 \(g_1\) 的首项整除，但被 \(g_2\) 的首项整除，商为 1，故加到 \(q_2\) 上（现在 \(q_2 = x+1\)），被除式变为 \((x-y^2+y+1)-(1)(g_2) = -y^2+1\). 现在的首项是 \(-y^2\)，它既不被 \(\operatorname{LT}(g_1) = xy\) 也不被 \(\operatorname{LT}(g_2) = x\) 整除，故 \(-y^2\) 被加到余式 \(r\) 上（现在 \(r = -y^2\)），被除式简单地变为 1. 最后，1 既不被 \(\operatorname{LT}(g_1)\) 也不被 \(\operatorname{LT}(g_2)\) 整除，故被加到余式上（现在 \(r = -y^2+1\)），过程终止。结果是
\[
q_1 = -1, \quad q_2 = x+1, \quad r = -y^2+1
\]
且
\[
f = x^2+x-y^2+y = (-1)(xy+1)+(x+1)(x+y)+(-y^2+1) = q_1g_1+q_2g_2+r.
\]

（3）设 \(f = x^2+x-y^2+y\) 同前一例，但交换除式 \(g_1\) 与 \(g_2\)：\(g_1 = x+y\)，\(g_2 = xy+1\). 此时简单的计算给出 \(q_1 = x-y+1\)、\(q_2 = 0\)、\(r = 0\)，且
\[
f = x^2+x-y^2+y = (x-y+1)(x+y) = q_1g_1+q_2g_2+r,
\]
这表明商 \(q_i\) 与余式 \(r\) 一般不是唯一的，且依赖于除式 \(g_1, \ldots, g_m\) 的次序。
\end{example}

例 3 中的计算表明多项式 \(f = x^2+x-y^2+y\) 是理想 \(I = (x+y, xy+1)\) 的元素，因为此时得到的余式是 0（事实上 \(f\) 只是第一个生成元的倍数）。然而在例 2 中，同一个多项式被 \(xy+1\) 和 \(x+y\) 除时得到非零余式 \(-y^2+1\)，而从那个计算完全看不出 \(f\) 是 \(I\) 的元素。下一个定理表明，若我们使用理想 \(I\) 的一个 Gröbner 基，这些困难就不会出现：我们得到一个唯一的余式，可用它来判断一个多项式 \(f\) 是否为理想 \(I\) 的元素。

\begin{theorem}\label{thm:9.23}
在 \(R = F[x_1, \ldots, x_n]\) 上固定一个单项序，并设 \(\{g_1, \ldots, g_m\}\) 是 \(R\) 中非零理想 \(I\) 的一个 Gröbner 基。则

（1）每个多项式 \(f \in R\) 都可以唯一地写成
\[
f = f_I + r
\]
的形式，其中 \(f_I \in I\)，且“余式”\(r\) 中没有任何非零单项式项被 \(\operatorname{LT}(g_1), \ldots, \operatorname{LT}(g_m)\) 中任何一个整除。

（2）\(f_I\) 和 \(r\) 都可以用 \(g_1, \ldots, g_m\) 通过一般多项式除法计算得到，并且与这些多项式在除法中使用的次序无关。

（3）余式 \(r\) 给出 \(f\) 在商环 \(F[x_1, \ldots, x_n]/I\) 中的陪集的唯一代表元。特别地，\(f \in I\) 当且仅当 \(r = 0\).
\end{theorem}

\begin{proof}
在用 \(g_1, \ldots, g_m\) 作 \(f\) 的一般多项式除法中，取 \(f_I = \sum_{i=1}^m q_ig_i \in I\)，立即得到对任意生成元 \(g_1, \ldots, g_m\) 的分解 \(f = f_I+r\). 现在假设 \(\{g_1, \ldots, g_m\}\) 是 Gröbner 基，且 \(f = f_I+r = f_I'+r'\). 则 \(r-r' = f_I'-f_I \in I\)，故它的首项 \(\operatorname{LT}(r-r')\) 是 \(\operatorname{LT}(I)\) 的元素；由于 \(\{g_1, \ldots, g_m\}\) 是 \(I\) 的 Gröbner 基，\(\operatorname{LT}(I)\) 就是理想 \((\operatorname{LT}(g_1), \ldots, \operatorname{LT}(g_m))\). 这个理想中每个元素都是单项式项 \(\operatorname{LT}(g_1), \ldots, \operatorname{LT}(g_m)\) 的倍数的和，故是若干个项之和，其中每一项都被某个 \(\operatorname{LT}(g_i)\) 整除。但 \(r\) 和 \(r'\)（从而 \(r-r'\)）都是单项式项之和，其中没有任何一项被 \(\operatorname{LT}(g_1), \ldots, \operatorname{LT}(g_m)\) 整除，除非 \(r-r' = 0\) 否则矛盾。由此 \(r = r'\) 唯一，从而 \(f_I = f-r\) 也唯一，这证明了（1）。

我们已经看到 \(f_I\) 与 \(r\) 可以用多项式除法算法地计算出来，而（1）中的唯一性蕴涵 \(r\) 与多项式 \(g_1, \ldots, g_m\) 在除法中使用的次序无关。类似地，\(f_I = \sum_{i=1}^m q_ig_i\) 也唯一确定（尽管各个商 \(q_i\) 一般并不唯一），这给出（2）。

（3）中的第一个断言由（1）中的唯一性立即得到。若 \(r = 0\)，则 \(f = f_I \in I\). 反之，若 \(f \in I\)，则由 \(f = f+0\) 以及 \(r\) 的唯一性可得 \(r = 0\)，定理的最后一个断言得证。
\end{proof}

如前面提到的，定理 \ref{thm:9.23} 的重要性以及 Gröbner 基的主要用途之一，就是代表元 \(r\) 的唯一性，它使得我们可以在商环 \(F[x_1, \ldots, x_n]/I\) 中进行有效的计算。

接下来我们证明：理想中一组多项式，若其首项生成该理想的所有首项，则它们实际上就是该理想的一组生成元（因而是 Gröbner 基——在某些著作中这正是 Gröbner 基的定义），这特别表明 Gröbner 基总是存在。

\begin{proposition}\label{pro:9.24}
在 \(R = F[x_1, \ldots, x_n]\) 上固定一个单项序，\(I\) 是 \(R\) 中的非零理想。

（1）若 \(g_1, \ldots, g_m\) 是 \(I\) 中任意满足 \(\operatorname{LT}(I) = (\operatorname{LT}(g_1), \ldots, \operatorname{LT}(g_m))\) 的元素，则 \(\{g_1, \ldots, g_m\}\) 是 \(I\) 的 Gröbner 基。

（2）理想 \(I\) 有 Gröbner 基。
\end{proposition}

\begin{proof}
假设 \(g_1, \ldots, g_m \in I\) 且 \(\operatorname{LT}(I) = (\operatorname{LT}(g_1), \ldots, \operatorname{LT}(g_m))\). 我们需要证明 \(g_1, \ldots, g_m\) 生成理想 \(I\). 若 \(f \in I\)，用一般多项式除法写成 \(f = \sum_{i=1}^m q_ig_i + r\)，其中余式 \(r\) 中没有任何非零项被任何 \(\operatorname{LT}(g_i)\) 整除。由于 \(f \in I\)，也有 \(r \in I\)，这意味着 \(\operatorname{LT}(r) \in \operatorname{LT}(I)\). 但那样 \(\operatorname{LT}(r)\) 就会被 \(\operatorname{LT}(g_1), \ldots, \operatorname{LT}(g_m)\) 中某一个整除，除非 \(r = 0\) 否则矛盾。于是 \(f = \sum_{i=1}^m q_ig_i\)，\(g_1, \ldots, g_m\) 生成 \(I\)，从而是 \(I\) 的 Gröbner 基，这证明了（1）。

对于（2），注意对任何理想 \(I\)，其首项理想 \(\operatorname{LT}(I)\) 是由 \(I\) 中所有多项式的首项生成的单项式理想。由习题 1，其中有限多个首项就足以生成 \(\operatorname{LT}(I)\)，比如说对某些 \(h_1, \ldots, h_k \in I\) 有 \(\operatorname{LT}(I) = (\operatorname{LT}(h_1), \ldots, \operatorname{LT}(h_k))\). 由（1），多项式 \(h_1, \ldots, h_k\) 是 \(I\) 的 Gröbner 基，完成证明。
\end{proof}

命题 \ref{pro:9.24} 证明了 Gröbner 基总是存在。接下来我们证明一个判别法，用它来判断理想 \(I\) 的一组给定生成元是否为 Gröbner 基，然后利用它给出求 Gröbner 基的算法。基本想法非常简单：正如我们在本节第一个例子中取 \(yf_1-xf_2\) 时看到的，\(\operatorname{LT}(I)\) 中会出现新的元素，它们来自生成元的使首项相消的线性组合。我们将看到，以这种简单方式从生成元得到新的首项，是一组生成元成为 Gröbner 基的唯一障碍。

一般地，若 \(f_1, f_2\) 是 \(F[x_1, \ldots, x_n]\) 中两个多项式，\(M\) 是单项式项 \(\operatorname{LT}(f_1)\) 与 \(\operatorname{LT}(f_2)\) 的首一最小公倍，则我们可以取差
\begin{equation}
S(f_1, f_2) = \frac{M}{\operatorname{LT}(f_1)}f_1 - \frac{M}{\operatorname{LT}(f_2)}f_2 \tag{9.1}
\end{equation}
来消去首项。下面的引理表明，这些初等线性组合覆盖了同多重次数多项式首项相消的所有情形。

\begin{lemma}\label{lem:9.25}
假设 \(f_1, \ldots, f_m \in F[x_1, \ldots, x_n]\) 是有相同多重次数 \(\alpha\) 的多项式，且常数 \(a_i \in F\) 的线性组合 \(h = a_1f_1+\cdots+a_mf_m\) 有严格更小的多重次数。则
\[
h = \sum_{i=2}^m b_iS(f_{i-1}, f_i), \quad \text{其中 } b_i \in F.
\]
\end{lemma}

\begin{proof}
写 \(f_i = c_if_i'\)，其中 \(c_i \in F\)、\(f_i'\) 是多重次数为 \(\alpha\) 的首一多项式。我们有
\[
h = \sum a_ic_if_i' = a_1c_1(f_1'-f_2') + (a_1c_1+a_2c_2)(f_2'-f_3') + \cdots + (a_1c_1+\cdots+a_{m-1}c_{m-1})(f_{m-1}'-f_m') + (a_1c_1+\cdots+a_mc_m)f_m'.
\]
注意 \(f_{i-1}'-f_i' = S(f_{i-1}, f_i)\). 由于 \(h\) 与每个 \(f_{i-1}'-f_i'\) 的多重次数都严格小于 \(\alpha\)，我们有 \(a_1c_1+\cdots+a_mc_m = 0\)，故右端最后一项为 0，引理得证。
\end{proof}

下一个命题表明：若差 \(S(g_i, g_j)\) 中没有 \(g_i\) 尚未覆盖的新首项，则一组生成元 \(g_1, \ldots, g_m\) 就是 Gröbner 基。这个结果给出了构造 Gröbner 基的算法的主要成分。

在 \(R = F[x_1, \ldots, x_n]\) 上固定一个单项序，并设 \(G = \{g_1, \ldots, g_m\}\) 是 \(R\) 中的有序多项式集合。对 \(f \in R\)，记 \(f \bmod G = r\)，若 \(r\) 是用 \(g_1, \ldots, g_m\)（按此次序）对 \(f\) 作一般多项式除法所得的余式。

\begin{proposition}\label{pro:9.26}
（\textbf{Buchberger 判别法}）设 \(R = F[x_1, \ldots, x_n]\)，在 \(R\) 上固定一个单项序。若 \(I = (g_1, \ldots, g_m)\) 是 \(R\) 中的非零理想，则 \(G = \{g_1, \ldots, g_m\}\) 是 \(I\) 的 Gröbner 基当且仅当对 \(1 \leq i < j \leq m\) 有 \(S(g_i, g_j) = 0 \bmod G\).
\end{proposition}

\begin{proof}
若 \(\{g_1, \ldots, g_m\}\) 是 \(I\) 的 Gröbner 基，则由定理 \ref{thm:9.23}，由于每个 \(S(g_i, g_j)\) 都是 \(I\) 的元素，故 \(S(g_i, g_j) = 0 \bmod G\).

现在假设对 \(1 \leq i < j \leq m\) 有 \(S(g_i, g_j) = 0 \bmod G\)，并取任意元素 \(f \in I\). 为看出 \(G\) 是 Gröbner 基，我们需要看出 \((\operatorname{LT}(g_1), \ldots, \operatorname{LT}(g_m))\) 包含 \(\operatorname{LT}(f)\). 由于 \(f \in I\)，我们可以对某些多项式 \(h_1, \ldots, h_m\) 写 \(f = \sum_{i=1}^m h_ig_i\). 这样的表示不唯一。在所有这些表示中选取使各个被加项（即 \(\max_{i=1,\ldots,m}\alpha(h_ig_i)\)）的最大多重次数最小者，记为 \(\alpha\). 显然 \(f\) 的多重次数不劣于所有被加项 \(h_ig_i\) 的最大多重次数，故 \(\alpha(f) \leq \alpha\). 写
\begin{equation}
\sum_{i=1}^m h_ig_i = \sum_{\alpha(h_ig_i)=\alpha} h_ig_i + \sum_{\alpha(h_ig_i)<\alpha} h_ig_i = \sum_{\alpha(h_ig_i)=\alpha} \operatorname{LT}(h_i)g_i + \sum_{\alpha(h_ig_i)=\alpha} (h_i-\operatorname{LT}(h_i))g_i + \sum_{\alpha(h_ig_i)<\alpha} h_ig_i. \tag{9.2}
\end{equation}
假设 \(\alpha(f) < \alpha\). 由于后两个和的多重次数也严格小于 \(\alpha\)，可知第一个和的多重次数严格小于 \(\alpha\). 若 \(a_i \in F\) 表示单项式项 \(\operatorname{LT}(h_i)\) 的常数系数，则 \(\operatorname{LT}(h_i) = a_ih_i''\)，其中 \(h_i''\) 是单项式。我们可以对 \(\sum a_i(h_i''g_i)\) 应用引理 \ref{lem:9.25}，把上面的第一个和写成 \(\sum b_iS(h_{i-1}''g_{i-1}, h_i''g_i)\)，其中 \(\alpha(h_{i-1}''g_{i-1}) = \alpha(h_i''g_i) = \alpha\). 设 \(\beta_{i-1,i}\) 是 \(\operatorname{LT}(g_{i-1})\) 与 \(\operatorname{LT}(g_i)\) 的首一最小公倍的多重次数。则简单计算表明 \(S(h_{i-1}''g_{i-1}, h_i''g_i)\) 就是把 \(S(g_{i-1}, g_i)\) 乘以多重次数为 \(\alpha-\beta_{i-1,i}\) 的单项式。多项式 \(S(g_{i-1}, g_i)\) 的多重次数小于 \(\beta_{i-1,i}\)，且由假设 \(S(g_{i-1}, g_i) = 0 \bmod G\). 这意味着把 \(S(g_{i-1}, g_i)\) 用 \(g_1, \ldots, g_m\) 作一般多项式除法后，每个 \(S(g_{i-1}, g_i)\) 都可以写成和 \(\sum q_jg_j\)，其中 \(\alpha(q_jg_j) < \beta_{i-1,i}\). 由此每个 \(S(h_{i-1}''g_{i-1}, h_i''g_i)\) 都是 \(\sum q_jg_j\) 形式的和，其中 \(\alpha(q_jg_j) < \alpha\). 但那样方程（9.2）右端所有的和都可写成 \(p_ig_i\) 形式的项之和，其中多项式 \(p_i\) 满足 \(\alpha(p_ig_i) < \alpha\). 这与 \(\alpha\) 的最小性矛盾，表明事实上 \(\alpha(f) = \alpha\)，即 \(f\) 的首项的多重次数为 \(\alpha\).

若现在我们取方程（9.2）中多重次数为 \(\alpha\) 的项，就得到
\[
\operatorname{LT}(f) = \sum_{\alpha(h_ig_i)=\alpha} \operatorname{LT}(h_i)\operatorname{LT}(g_i),
\]
故确实 \(\operatorname{LT}(f) \in (\operatorname{LT}(g_1), \ldots, \operatorname{LT}(g_m))\). 由此 \(G = \{g_1, \ldots, g_m\}\) 是 Gröbner 基。
\end{proof}

\noindent\textbf{Buchberger 算法}

Buchberger 判别法可以用来给出求理想 \(I\) 的 Gröbner 基的算法，如下。若 \(I = (g_1, \ldots, g_m)\)，且用 \(G = \{g_1, \ldots, g_m\}\) 通过一般多项式除法作除法时每个 \(S(g_i, g_j)\) 的余式都是 0，则 \(G\) 是 Gröbner 基。否则 \(S(g_i, g_j)\) 有一个非零余式 \(r\). 把多项式 \(g_{m+1} = r\) 添加到 \(G\) 上：\(G' = \{g_1, \ldots, g_m, g_{m+1}\}\)，然后重新开始（注意这仍然是 \(I\) 的一组生成元，因为 \(g_{m+1} \in I\)）。不难验证这个过程在有限步后终止于一个满足 Buchberger 判别法的生成集 \(G\)，因而是 \(I\) 的 Gröbner 基（参见习题 16）。注意一旦某个 \(S(g_i, g_j)\) 用 \(G\) 中的多项式作除法得到余式 0，那么当更多的多项式被添加到 \(G\) 上时它仍然得到余式 0.

若 \(\{g_1, \ldots, g_m\}\) 是理想 \(I\) 的 Gröbner 基，且对某个 \(j \neq i\) 有 \(\operatorname{LT}(g_j)\) 被 \(\operatorname{LT}(g_i)\) 整除，则作为 \(\operatorname{LT}(I)\) 的生成元不需要 \(\operatorname{LT}(g_j)\). 由命题 \ref{pro:9.24}，我们可以删去 \(g_j\) 而仍保留 \(I\) 的一个 Gröbner 基。我们还可以不失一般性地假设每个 \(g_i\) 的首项是首一的。\(I\) 的一个 Gröbner 基 \(\{g_1, \ldots, g_m\}\)，若每个 \(\operatorname{LT}(g_i)\) 都是首一的，且当 \(i \neq j\) 时 \(\operatorname{LT}(g_j)\) 不被 \(\operatorname{LT}(g_i)\) 整除，则称为\textbf{极小 Gröbner 基}。虽然极小 Gröbner 基不唯一，但元素的个数与它们的首项是唯一的（参见习题 15）。

\begin{example}
（1）选择 \(F[x, y]\) 上的字典序 \(x > y\)，并考虑由 \(f_1 = x^3y-xy^2+1\) 与 \(f_2 = x^2y^2-y^3-1\) 生成的理想 \(I\)，如本节开头例 1 中那样。为检验 \(G = \{f_1, f_2\}\) 是否为 Gröbner 基，我们计算 \(S(f_1, f_2) = yf_1-xf_2 = x+y\)，它被 \(\{f_1, f_2\}\) 除时是自身的余式，故 \(G\) 不是 \(I\) 的 Gröbner 基。令 \(f_3 = x+y\)，并扩大生成集：\(G' = \{f_1, f_2, f_3\}\). 现在 \(S(f_1, f_2) = 0 \bmod G'\)，而简单的计算给出
\[
S(f_1, f_3) = f_1 - x^2yf_3 = -x^2y^2-x^2+1 \equiv 0 \bmod G', \qquad S(f_2, f_3) = f_2 - x^2yf_3 = -xy^3-y-1 \equiv y^4-y^3-1 \bmod G'.
\]
令 \(f_4 = y^4-y^3-1\)，并把生成集扩大为 \(G'' = \{f_1, f_2, f_3, f_4\}\). 先前的 0 余式仍是 0，而现在由 \(f_4\) 的选取有 \(S(f_2, f_3) = 0 \bmod G''\). 一些额外的计算给出
\[
S(f_1, f_4) = S(f_2, f_4) = S(f_3, f_4) = 0 \bmod G'',
\]
故 \(\{x^3y-xy^2+1,\ x^2y^2-y^3-1,\ x+y,\ y^4-y^3-1\}\) 是 \(I\) 的 Gröbner 基。特别地，\(\operatorname{LT}(I)\) 由这四个多项式的首项生成，故 \(\operatorname{LT}(I) = (x^3y, x^2y^2, x, y^4) = (x, y^4)\)，正如前面提到的。于是 \(I\) 中 \(x+y\) 与 \(y^4-y^3-1\) 的首项生成 \(\operatorname{LT}(I)\)，故由命题 \ref{pro:9.24}，\(\{x+y,\ y^4-y^3-1\}\) 给出 \(I\) 的一个极小 Gröbner 基：
\[
I = (x+y,\ y^4-y^3-1).
\]
这个对 \(I\) 的描述比 \(I = (x^3y-xy^2+1,\ x^2y^2-y^3-1)\) 简单得多。

（2）选择 \(F[x, y]\) 上的字典序 \(y > x\)，并考虑上一例中的理想 \(I\). 此时 \(S(f_1, f_2)\) 给出余式 \(f_3 = -x-y\)；然后 \(S(f_1, f_3)\) 给出余式 \(f_4 = -x^4-x^3+1\)，此后所有余式相对于 Gröbner 基 \(\{x^3y-xy^2+1,\ x^2y^2-y^3-1,\ -x-y,\ -x^4-x^3+1\}\) 都是 0. 这里 \(\operatorname{LT}(I) = (-xy^2, -y^3, -y, -x^4) = (y, x^4)\)，正如前面提到的，而 \(\{x+y,\ x^4+x^3-1\}\) 给出 \(I\) 关于这个次序的一个极小 Gröbner 基：
\[
I = (x+y,\ x^4+x^3-1),
\]
这是对 \(I\) 的另一个更简单的描述。
\end{example}

在上面例 1 中容易验证 \(\{x+y^4-y^3+y-1,\ y^4-y^3-1\}\) 也是 \(I\) 的极小 Gröbner 基（它就是 \(\{f_3+f_4, f_4\}\)），故即使固定了 \(F[x_1, \ldots, x_n]\) 上的单项序，理想 \(I\) 的极小 Gröbner 基也不唯一。通过加强基中首项的可整除性条件，我们可以得到一个重要的唯一性性质。

\begin{definition}
在 \(R = F[x_1, \ldots, x_n]\) 上固定一个单项序。\(R\) 中非零理想 \(I\) 的 Gröbner 基 \(\{g_1, \ldots, g_m\}\) 称为\textbf{约化 Gröbner 基}，若

（a）每个 \(g_i\) 的首项是首一的，即 \(\operatorname{LT}(g_i)\) 是首一的，\(i = 1, \ldots, m\)，且

（b）当 \(j \neq i\) 时 \(g_j\) 中没有任何项被 \(\operatorname{LT}(g_i)\) 整除。
\end{definition}

注意约化 Gröbner 基特别地是极小 Gröbner 基。若 \(G = \{g_1, \ldots, g_m\}\) 是 \(I\) 的极小 Gröbner 基，则对任意 \(i \neq j\)，首项 \(\operatorname{LT}(g_j)\) 都不被 \(\operatorname{LT}(g_i)\) 整除。因此，若我们用多项式除法把 \(g_j\) 用 \(G\) 中其余多项式除，就得到 \(I\) 中一个与原 \(g_j\) 有相同首项的余式 \(\tilde{g}_j\)（由定理 \ref{thm:9.23} 的（2），余式 \(\tilde{g}_j\) 与除法中多项式的次序无关）。由命题 \ref{pro:9.24}，在 \(G\) 中把 \(g_j\) 换成 \(\tilde{g}_j\) 仍给出 \(I\) 的极小 Gröbner 基，并且在这个基中 \(\tilde{g}_j\) 的任何项都不被 \(\operatorname{LT}(g_i)\)（\(i \neq j\)）整除。因此把 \(G\) 的每个元素换成它被 \(G\) 中其余元素除后的余式，就得到 \(I\) 的一个约化 Gröbner 基。约化 Gröbner 基的重要性在于它们（对给定的单项序）是唯一的，如下一个结果所示。

\begin{theorem}\label{thm:9.27}
在 \(R = F[x_1, \ldots, x_n]\) 上固定一个单项序。则 \(R\) 中每个非零理想 \(I\) 都有唯一的约化 Gröbner 基。
\end{theorem}

\begin{proof}
由习题 15，两个约化基有相同个数的元素和相同的首项，因为约化基也是极小基。若 \(G = \{g_1, \ldots, g_m\}\) 与 \(G' = \{g_1', \ldots, g_m'\}\) 是同一个非零理想 \(I\) 的两个约化基，则必要时重新排列后我们可以假设 \(\operatorname{LT}(g_i) = \operatorname{LT}(g_i') = h_i\)（\(i = 1, \ldots, m\)）。任取 \(i\)，考虑多项式 \(f_i = g_i - g_i'\). 若 \(f_i\) 非零，则由于 \(f_i \in I\)，它的首项必被某个 \(h_j\) 整除。由约化基的定义，当 \(j \neq i\) 时 \(h_j\) 不整除 \(g_i\) 或 \(g_i'\) 中的任何项，故不整除 \(\operatorname{LT}(f_i)\). 但 \(h_i\) 也不整除 \(\operatorname{LT}(f_i)\)，因为 \(f_i\) 的所有项都有严格更小的多重次数。这迫使 \(f_i = 0\)，即对每个 \(i\) 都有 \(g_i = g_i'\)，故 \(G = G'\).
\end{proof}

约化 Gröbner 基唯一性的一个应用是：给出一个判断多项式环中两个理想何时相等的计算方法。

\begin{corollary}\label{cor:9.28}
设 \(I\) 和 \(J\) 是 \(F[x_1, \ldots, x_n]\) 中两个理想。则 \(I = J\) 当且仅当 \(I\) 与 \(J\) 关于 \(F[x_1, \ldots, x_n]\) 上任意固定的单项序有相同的约化 Gröbner 基。
\end{corollary}

\begin{example}
（1）考虑 \(F[x, y]\) 中的理想 \(I = (h_1, h_2, h_3)\)，其中 \(h_1 = x^2+xy^5+y^4\)、\(h_2 = xy^6-xy^3+y^5-y^2\)、\(h_3 = xy^5-xy^2\). 使用字典序 \(x > y\)，我们得到
\[
S(h_1, h_2) = S(h_1, h_3) = 0 \bmod \{h_1, h_2, h_3\} \quad \text{且} \quad S(h_2, h_3) = y^5-y^2 \bmod \{h_1, h_2, h_3\}.
\]
令 \(h_4 = y^5-y^2\)，我们发现对 \(1 \leq i < j \leq 4\) 有 \(S(h_i, h_j) = 0 \bmod \{h_1, h_2, h_3, h_4\}\)，故
\[
\{x^2+xy^5+y^4,\ xy^6-xy^3+y^5-y^2,\ xy^5-xy^2,\ y^5-y^2\}
\]
是 \(I\) 的 Gröbner 基。这个基的首项是 \(x^2, xy^6, xy^5, y^5\). 由于 \(y^5\) 整除 \(xy^5\) 和 \(xy^6\)，我们可以删去第二和第三个生成元，得到 \(I\) 的一个极小 Gröbner 基 \(\{x^2+xy^5+y^4,\ y^5-y^2\}\). 第一个生成元中的第二项被第二个生成元的首项 \(y^5\) 整除，故这不是约化 Gröbner 基。把 \(x^2+xy^5+y^4\) 换成它被基中其余多项式（此时只有多项式 \(y^5-y^2\)）除后的余式 \(x^2+xy^2+y^4\)，我们得到 \(I\) 的约化 Gröbner 基 \(\{x^2+xy^2+y^4,\ y^5-y^2\}\).

（2）考虑 \(F[x, y]\) 中的理想 \(J = (h_1, h_2, h_3)\)，其中 \(h_1 = xy^3+y^3+1\)、\(h_2 = x^3y-x^3+1\)、\(h_3 = x+y\). 使用字典序 \(x > y\)，我们得到 \(S(h_1, h_2) = 0 \bmod \{h_1, h_2, h_3\}\) 以及 \(S(h_1, h_3) = y^4-y^3-1 \bmod \{h_1, h_2, h_3\}\). 令 \(h_4 = y^4-y^3-1\)，我们发现对 \(1 \leq i < j \leq 4\) 有 \(S(h_i, h_j) = 0 \bmod \{h_1, h_2, h_3, h_4\}\)，故
\[
\{xy^3+y^3+1,\ x^3y-x^3+1,\ x+y,\ y^4-y^3-1\}
\]
是 \(J\) 的 Gröbner 基。这个基的首项是 \(xy^3, x^3y, x, y^4\)，故 \(\{x+y,\ y^4-y^3-1\}\) 是 \(J\) 的极小 Gröbner 基。此时 \(y^4-y^3-1\) 中没有任何项被 \(x+y\) 的首项整除，\(x+y\) 中也没有任何项被 \(y^4-y^3-1\) 的首项整除，故 \(\{x+y,\ y^4-y^3-1\}\) 是 \(J\) 的约化 Gröbner 基。这正是命题 \ref{pro:9.26} 之后的例 1 中理想 \(I\) 的基，故这两个理想相等：
\[
(x^3y-xy^2+1,\ x^2y^2-y^3-1) = (xy^3+y^3+1,\ x^3y-x^3+1,\ x+y)
\]
（并且它们都等于理想 \((x+y,\ y^4-y^3-1)\)）。
\end{example}

\noindent\textbf{Gröbner 基与解代数方程：消元}

Gröbner 基的理论在显式求解代数方程组时非常有用，也是计算机代数程序求解方程组的基础。假设 \(S = \{f_1, \ldots, f_m\}\) 是 \(n\) 个变量 \(x_1, \ldots, x_n\) 中一组多项式，我们想求方程组 \(f_1 = 0, f_2 = 0, \ldots, f_m = 0\) 的解（即 \(S\) 中多项式的公共零点集）。若 \((a_1, \ldots, a_n)\) 是这个方程组的任一解，则 \(S\) 所生成理想 \(I\) 中的每个元素 \(f\) 也满足 \(f(a_1, \ldots, a_n) = 0\). 此外，容易验证若 \(S' = \{g_1, \ldots, g_s\}\) 是理想 \(I\) 的任意一组生成元，则方程组 \(g_1 = 0, \ldots, g_s = 0\) 的解集与原来的解集相同。

在 \(f_1, \ldots, f_m\) 都是一次多项式的情形，可以用初等方法依次消去变量 \(x_1, x_2, \ldots\) 得到方程组的解——用原方程的线性组合消去变量 \(x_1\)，再用这些方程消去 \(x_2\)，等等，得到一个容易求解的方程组（这就是线性代数中的“高斯--若尔当消元法”，参见第 11.2 节的习题）。

非线性多项式方程的情形自然更复杂，但基本原则是相同的。若理想 \(I\) 中有一个只涉及某一个变量（比如 \(x_n\)）的非零多项式 \(p(x_n)\)，则最后一个坐标 \(a_n\) 是 \(p(x_n) = 0\) 的解。若理想 \(I\) 中有一个只涉及 \(x_{n-1}\) 和 \(x_n\) 的多项式 \(q(x_{n-1}, x_n)\)，则坐标 \(a_{n-1}\) 是 \(q(x_{n-1}, a_n) = 0\) 的解，等等。如果我们能依次找到 \(I\) 中消去变量 \(x_1, x_2, \ldots\) 的多项式，我们就能显式地确定原方程组的所有解 \((a_1, \ldots, a_n)\).

寻找由 \(S\) 中方程组推出的方程，即寻找 \(I\) 中不涉及某些变量的元素，称为\textbf{消元理论}。\(I\) 中不涉及变量 \(x_1, \ldots, x_i\) 的多项式，即 \(I \cap F[x_{i+1}, \ldots, x_n]\)，容易看出是 \(F[x_{i+1}, \ldots, x_n]\) 中的一个理想，它有一个名称。

\begin{definition}
若 \(I\) 是 \(F[x_1, \ldots, x_n]\) 中的理想，则 \(I_i = I \cap F[x_{i+1}, \ldots, x_n]\) 称为 \(I\) 关于次序 \(x_1 > \cdots > x_n\) 的\textbf{第 \(i\) 个消元理想}。
\end{definition}

用消元法求解方程组能否成功，取决于能否确定这些消元理想（最终取决于这些消元理想是否非零）。

下面这个基本命题表明，若使用字典单项序 \(x_1 > \cdots > x_n\) 计算 \(I\) 的 Gröbner 基，则所得基中不涉及变量 \(x_1, \ldots, x_i\) 的元素不仅确定了第 \(i\) 个消元理想，事实上还给出 \(I\) 的第 \(i\) 个消元理想的 Gröbner 基。

\begin{proposition}\label{pro:9.29}
（\textbf{消元定理}）假设 \(G = \{g_1, \ldots, g_m\}\) 是 \(F[x_1, \ldots, x_n]\) 中非零理想 \(I\) 关于字典单项序 \(x_1 > \cdots > x_n\) 的 Gröbner 基。则 \(G \cap F[x_{i+1}, \ldots, x_n]\) 是 \(I\) 的第 \(i\) 个消元理想 \(I_i = I \cap F[x_{i+1}, \ldots, x_n]\) 的 Gröbner 基。特别地，\(I \cap F[x_{i+1}, \ldots, x_n] = 0\) 当且仅当 \(G \cap F[x_{i+1}, \ldots, x_n] = 0\).
\end{proposition}

\begin{proof}
记 \(G_i = G \cap F[x_{i+1}, \ldots, x_n]\). 则 \(G_i \subseteq I_i\)，故由命题 \ref{pro:9.24}，为看出 \(G_i\) 是 \(I_i\) 的 Gröbner 基，只需看出 \(G_i\) 中元素的首项 \(\operatorname{LT}(G_i)\) 作为 \(F[x_{i+1}, \ldots, x_n]\) 中的理想生成 \(\operatorname{LT}(I_i)\). 当然作为 \(F[x_{i+1}, \ldots, x_n]\) 中的理想有 \((\operatorname{LT}(G_i)) \subseteq \operatorname{LT}(I_i)\). 为证明反向包含，取 \(I_i\) 中任意元素 \(f\). 则 \(f \in I\)，且由于 \(G\) 是 \(I\) 的 Gröbner 基，我们有
\[
\operatorname{LT}(f) = a_1(x_1, \ldots, x_n)\operatorname{LT}(g_1) + \cdots + a_m(x_1, \ldots, x_n)\operatorname{LT}(g_m)
\]
对某些多项式 \(a_1, \ldots, a_m \in F[x_1, \ldots, x_n]\). 把每个多项式 \(a_i\) 写成单项式项之和，我们看到 \(\operatorname{LT}(f)\) 是形如 \(ax_1^{e_1}\cdots x_n^{e_n}\operatorname{LT}(g_i)\) 的单项式项之和。由于 \(\operatorname{LT}(f)\) 只涉及变量 \(x_{i+1}, \ldots, x_n\)，所有含变量 \(x_1, \ldots, x_i\) 的这种项之和必为 0，故 \(\operatorname{LT}(f)\) 也是那些只涉及 \(x_{i+1}, \ldots, x_n\) 的单项式项之和。由此 \(\operatorname{LT}(f)\) 可以写成 \(\operatorname{LT}(g_1)\) 的某些单项式项（其中 \(\operatorname{LT}(g_1)\) 不涉及变量 \(x_1, \ldots, x_i\)）的 \(F[x_{i+1}, \ldots, x_n]\)-线性组合。但由次序的选取，若 \(\operatorname{LT}(g_1)\) 不涉及 \(x_1, \ldots, x_i\)，则 \(g_1\) 的其余任何项也不涉及它们，即 \(g_1 \in G_i\). 于是 \(\operatorname{LT}(f)\) 可以写成 \(\operatorname{LT}(G_i)\) 中元素的 \(F[x_{i+1}, \ldots, x_n]\)-线性组合，完成证明。
\end{proof}

还要注意，只要使用适当的单项序，Gröbner 基就可以用来消去任意变量。

\begin{example}
（1）椭圆 \(2x^2+2xy+y^2-2x-2y = 0\) 与圆 \(x^2+y^2 = 1\) 交于两点。为求它们，我们用字典序 \(x > y\) 计算理想 \(I = (2x^2+2xy+y^2-2x-2y,\ x^2+y^2-1)\) 在 \(\mathbb{R}[x, y]\) 中的 Gröbner 基来消去 \(x\)，得到 \(g_1 = 2x+y^2+5y^3-2\) 和 \(g_2 = 5y^4-4y^3\). 于是 \(y^3(5y-4) = 0\)，故 \(y = 0\) 或 \(y = 4/5\). 把这些值代入 \(g_1 = 0\) 并解出 \(x\)，我们得到两个交点是 \((1, 0)\) 和 \((-3/5, 4/5)\).

若改用字典序 \(y > x\) 来消去 \(y\)，则得到的 Gröbner 基是 \(\{y^2+x^2-1,\ 2yx-2y+x^2-2x+1,\ 5x^3-7x^2-x+3\}\). 此时 \(5x^3-7x^2-x+3 = (x-1)^2(5x+3)\) 表明 \(x\) 为 1 或 \(-3/5\)，我们得到与前面相同的解，只是花费了更多功夫。

（2）上一例中的解也可以用初等方法求出。现在考虑方程组
\[
x^3-2xy+y^3 = 0 \quad \text{和} \quad x^5-2x^2y^2+y^5 = 0
\]
在 \(\mathbb{C}\) 中的解。关于字典序 \(x > y\) 计算由 \(f_1 = x^3-2xy+y^3\) 和 \(f_2 = x^5-2x^2y^2+y^5\) 生成的理想的 Gröbner 基，我们得到基
\[
\begin{aligned}
g_1 &= x^3-2xy+y^3 \\
g_2 &= 200xy^2+193y^9+158y^8-45y^7-456y^6+50y^5-100y^4 \\
g_3 &= y^{10}-y^8-2y^7+2y^6.
\end{aligned}
\]
原方程组的任何解都满足 \(g_1 = g_2 = g_3 = 0\). 由于 \(g_3 = y^6(y-1)^2(y^2+2y+2)\)，我们有 \(y = 0\)、\(y = 1\) 或 \(y = -1\pm i\). 由于 \(g_1(x, 0) = x^3\)，\(g_2(x, 0) = 0\)，我们看到 \((0,0)\) 是 \(y = 0\) 时的唯一解。由于 \(g_1(x, 1) = x^3-2x+1\) 与 \(g_2(x, 1) = 200(x-1)\) 只有公共零点 \(x = 1\)，\(y = 1\) 时的唯一解是 \((1,1)\). 最后，
\[
g_1(x, -1\pm i) = x^3 + (2\mp 2i)x + (2\pm 2i), \qquad g_2(x, -1\pm i) = -400i(x+1\pm i),
\]
而快速检验表明当 \(y = -1\pm i\) 时公共零点为 \(x = -1\mp i\). 因此原方程组恰有四个解，即 \((x, y) = (0,0)\)、\((1,1)\)、\((-1+i, -1-i)\) 或 \((-1-i, -1+i)\).

（3）考虑方程组
\[
x+y+z = 1, \qquad x^2+y^2+z^2 = 2, \qquad x^3+y^3+z^3 = 3
\]
在 \(\mathbb{C}\) 中的解。关于字典序 \(x > y > z\) 的约化 Gröbner 基是
\[
\{x+y+z-1,\ y^2+yz-y+z^2-z-(1/2),\ z^3-z^2-(1/2)z-(1/6)\},
\]
故 \(z\) 是多项式 \(t^3-t^2-(1/2)t-(1/6)\) 的根（由对称性，\(x\) 和 \(y\) 也是同一个多项式的根）。对这个多项式的三个根中的每一个，都有两个 \(y\) 的值和一个相应的 \(x\) 的值使 Gröbner 基中前两个多项式为 0. 快速检验表明得到的六个解正是多项式 \(t^3-t^2-(1/2)t-(1/6)\)（它在 \(\mathbb{Q}\) 上不可约）的三个互异根的某种排列。
\end{example}

如前面的例子所示，多项式方程组 \(f_1 = 0, f_2 = 0, \ldots, f_m = 0\) 的解的研究，与这些多项式在 \(F[x_1, \ldots, x_n]\) 中生成的理想 \(I = (f_1, f_2, \ldots, f_m)\) 的研究密切相关。这个基本联系是数学中一个重要而活跃的分支——“代数几何”——的出发点，它将在第 15 章中介绍，那里还会给出 Gröbner 基的更多应用。

我们以说明如何计算多项式环中理想的加法、乘法与交这三种基本集合运算来结束本节。假设 \(I = (f_1, \ldots, f_s)\) 和 \(J = (h_1, \ldots, h_t)\) 是 \(F[x_1, \ldots, x_n]\) 中两个理想。则 \(I+J = (f_1, \ldots, f_s, h_1, \ldots, h_t)\)，而 \(IJ = (f_1h_1, \ldots, f_ih_j, \ldots, f_sh_t)\). 下面这个命题说明如何计算任意两个理想的交。

\begin{proposition}\label{pro:9.30}
若 \(I\) 和 \(J\) 是 \(F[x_1, \ldots, x_n]\) 中任意两个理想，则 \(tI+(1-t)J\) 是 \(F[t, x_1, \ldots, x_n]\) 中的理想，且 \(I \cap J = (tI+(1-t)J) \cap F[x_1, \ldots, x_n]\). 特别地，\(I \cap J\) 是 \(tI+(1-t)J\) 关于次序 \(t > x_1 > \cdots > x_n\) 的第一个消元理想。
\end{proposition}

\begin{proof}
首先，\(tI\) 和 \((1-t)J\) 显然是 \(F[x_1, \ldots, x_n, t]\) 中的理想，故它们的和 \(tI+(1-t)J\) 也是 \(F[x_1, \ldots, x_n, t]\) 中的理想。若 \(f \in I \cap J\)，则 \(f = tf+(1-t)f\) 表明 \(I \cap J \subseteq (tI+(1-t)J) \cap F[x_1, \ldots, x_n]\). 反之，假设 \(f = tf_1+(1-t)h_1\) 是 \(F[x_1, \ldots, x_n]\) 的元素，其中 \(f_1 \in I\)、\(h_1 \in J\). 则 \(t(f_1-h_1) = f-h_1 \in F[x_1, \ldots, x_n]\) 表明 \(f_1-h_1 = 0\)，且 \(f = h_1 = f_1 \in I \cap J\).

由于 \(I \cap J = (tI+(1-t)J) \cap F[x_1, \ldots, x_n]\)，\(I \cap J\) 就是 \(tI+(1-t)J\) 关于次序 \(t > x_1 > \cdots > x_n\) 的第一个消元理想。
\end{proof}

若 \(I = (f_1, \ldots, f_s)\)、\(J = (h_1, \ldots, h_t)\)，则 \(tI+(1-t)J = (tf_1, \ldots, tf_s, (1-t)h_1, \ldots, (1-t)h_t)\). 由命题 \ref{pro:9.29}，在 \(F[t, x_1, \ldots, x_n]\) 中关于字典单项序 \(t > x_1 > \cdots > x_n\) 计算这个理想的 Gröbner 基时，其中不涉及 \(t\) 的元素就给出 \(F[x_1, \ldots, x_n]\) 中理想 \(I \cap J\) 的 Gröbner 基。

\begin{example}
设 \(I = (x, y) = (x^2, xy, y^2)\)，\(J = (x)\). 关于字典单项序 \(t > x > y\)，\(tI+(1-t)J\) 在 \(F[t, x, y]\) 中的约化 Gröbner 基是 \(\{tx-x, ty^2, x^2, xy\}\)，故 \(I \cap J = (x^2, xy)\).
\end{example}

\subsection*{习题}

\begin{enumerate}
\item 假设 \(I\) 是 \(F[x_1, \ldots, x_n]\) 中由多项式集合 \(S\)（可以无限）生成的理想。证明 \(S\) 中有限多个多项式足以生成 \(I\).〔用定理 \ref{thm:9.21} 写 \(I = (f_1, \ldots, f_m)\)，然后用 \(S\) 中的多项式表示每个 \(f_i \in I\).〕

\item 设 \(\geq\) 是任意单项序。

（a）证明对任意非零多项式 \(f\) 和 \(g\) 有 \(\operatorname{LT}(fg) = \operatorname{LT}(f)\operatorname{LT}(g)\) 且 \(\alpha(fg) = \alpha(f)+\alpha(g)\).

（b）证明 \(\alpha(f+g) \leq \max(\alpha(f), \alpha(g))\)，且当 \(\alpha(f) \neq \alpha(g)\) 时取等号。

（c）证明对每个单项式 \(m\) 有 \(m \geq 1\).

（d）证明若 \(m_1\) 整除 \(m_2\) 则 \(m_2 \geq m_1\). 由此推出多项式的首项不整除它的任何低次项。

\item 证明若 \(\geq\) 是非空集合上任意全序或偏序，则下述两个命题等价：

（i）每个非空子集都含有最小元。

（ii）不存在无穷严格递减序列 \(a_1 > a_2 > a_3 > \cdots\)（这称为\textbf{降链条件}或 D.C.C.）。

由此推出一般多项式除法总在有限步内终止。

\item 设 \(\geq\) 是单项序，对单项式 \(m_1, m_2\) 定义 \(m_1 \geq_g m_2\)，若或者 \(\deg m_1 > \deg m_2\)，或者 \(\deg m_1 = \deg m_2\) 且 \(m_1 \geq m_2\).

（a）证明 \(\geq_g\) 也是单项序。（关系 \(\geq_g\) 称为 \(\geq\) 的\textbf{分级}。比较时最重要的标准是次数的次序有时称为\textbf{分级序}或\textbf{次数序}，故这个习题给出构造分级序的方法。）

（b）字典序 \(x_1 > \cdots > x_n\) 的分级称为 \textbf{grlex 单项序}。证明关于 grlex 序有 \(x_1^3 > x_1^2x_2 > x_1x_2^2 > x_2^3 > x_1\)，而关于字典序有 \(x_1^3 > x_1^2x_2 > x_1x_2^2 > x_1 > x_2^3\).

\item \textbf{grevlex 单项序}的定义是：先选取变量的一个次序 \(\{x_1, x_2, \ldots, x_n\}\)，然后对单项式 \(m_1, m_2\) 定义 \(m_1 \geq m_2\)，若或者 \(\deg m_1 > \deg m_2\)，或者 \(\deg m_1 = \deg m_2\) 且按次序 \(x_n, x_{n-1}, \ldots, x_1\) 第一个使 \(m_1\) 与 \(m_2\) 不同的指数在 \(m_1\) 中更小。

（a）证明 grevlex 是单项序，且满足 \(x_1 > x_2 > \cdots > x_n\).

（b）证明关于 \(\{x_1, x_2\}\) 的 grevlex 序在 \(F[x_1, x_2]\) 上就是 \(x_1 > x_2\) 的分级字典序，但关于 \(\{x_1, x_2, x_3\}\) 的 grevlex 序在 \(F[x_1, x_2, x_3]\) 上不是任何字典序的分级。

（c）证明关于 \(\{x_1, x_2, x_3\}\) 的 grevlex 单项序有
\[
x_1^2x_2 > x_1^2x_3 > x_1x_2x_3 > x_2^2x_3 > x_1x_3^2 > x_2x_3^2 > x_1^2 > x_1 > x_2 > x_3.
\]

\item 证明关于字典单项序 \(x > y > z\) 有
\[
x^3y > x^3z^2 > x^3z > x^2y^2z > x^2y > xz^2 > y^2z^2 > y^2z.
\]
证明关于相应的 grlex 单项序有
\[
x^3z^2 > x^2y^2z > x^3y > x^3z > y^2z^2 > x^2y > xz^2 > y^2z,
\]
而关于 \(\{x, y, z\}\) 的 grevlex 单项序有
\[
x^2y^2z > x^3z^2 > x^3y > x^3z > y^2z^2 > x^2y > y^2z > xz^2.
\]

\item 把单项式 \(x^2z,\ xy^2z,\ x^2,\ x^2y^2z,\ x^3y,\ x^3z^2,\ x^2yz^2,\ x^2z^2\) 按字典单项序 \(x > y > z\) 排序，再按相应的 grlex 单项序排序，最后按关于 \(\{x, y, z\}\) 的 grevlex 单项序排序。

\item 证明 \(F[x_1, \ldots, x_n]\) 上有 \(n!\) 个不同的字典单项序。类似地证明有 \(n!\) 个不同的 grlex 单项序和 grevlex 单项序。

\item 可以证明 \(F[x_1, \ldots, x_n]\) 上任何单项序都可以如下得到。对 \(k \geq n\)，设 \(v_1, v_2, \ldots, v_k\) 是欧几里得 \(n\) 维空间 \(\mathbb{R}^n\) 中的非零向量，两两正交：对所有 \(i \neq j\) 有 \(v_i \cdot v_j = 0\)（这里 \(\cdot\) 是通常的点积），并设 \(v_1\) 的所有坐标非负。定义单项式上的次序 \(\geq\)：\(m_1 \geq m_2\) 当且仅当对某个 \(t \geq k\)，对所有 \(i \in \{1, 2, \ldots, t-1\}\) 有 \(v_i \cdot \alpha(m_1) = v_i \cdot \alpha(m_2)\)，且 \(v_t \cdot \alpha(m_1) > v_t \cdot \alpha(m_2)\).

（a）设 \(k = n\)，\(v_i = (0, \ldots, 0, 1, 0, \ldots, 0)\)（1 在第 \(i\) 个位置）。证明 \(\geq\) 定义的是字典序 \(x_1 > x_2 > \cdots > x_n\).

（b）设 \(k = n\)，并定义 \(v_1 = (1, 1, \ldots, 1)\) 及 \(v_i = (1, 1, \ldots, 1, -n+i-1, 0, \ldots, 0)\)，其中末尾有 \(i-2\) 个零，\(2 \leq i \leq n\). 证明 \(\geq\) 定义的是关于 \(\{x_1, \ldots, x_n\}\) 的 grlex 序。

\item 假设 \(I\) 是由单项式 \(m_1, \ldots, m_k\) 生成的单项式理想。证明多项式 \(f \in F[x_1, \ldots, x_n]\) 属于 \(I\) 当且仅当 \(f\) 的每个单项式项都是某个 \(m_i\) 的倍数。〔对多项式 \(a_1, \ldots, a_k \in F[x_1, \ldots, x_n]\)，展开 \(a_1m_1+\cdots+a_km_k\) 并注意每个单项式项都至少是某个 \(m_i\) 的倍数。〕证明 \(x^2yz+3xy^2\) 是理想 \(J = (xyz, y^2) \subset F[x, y, z]\) 的元素，但不是理想 \(I' = (xz^2, y^2)\) 的元素。

\item 在 \(R = F[x_1, \ldots, x_n]\) 上固定一个单项序，并假设 \(\{g_1, \ldots, g_m\}\) 是 \(R\) 中理想 \(I\) 的 Gröbner 基。证明 \(h \in \operatorname{LT}(I)\) 当且仅当 \(h\) 是若干单项式项之和，其中每一项都被某个 \(\operatorname{LT}(g_i)\)（\(1 \leq i \leq m\)）整除。〔利用上一习题。〕

\item 假设 \(I\) 是由单项式 \(g_1, \ldots, g_m\) 生成的单项式理想。用上一习题直接证明 \(\{g_1, \ldots, g_m\}\) 是 \(I\) 的 Gröbner 基。

\item 假设 \(I\) 是由单项式 \(g_1, \ldots, g_m\) 生成的单项式理想。用 Buchberger 判别法证明 \(\{g_1, \ldots, g_m\}\) 是 \(I\) 的 Gröbner 基。

\item 假设 \(I\) 是 \(R = F[x_1, \ldots, x_n]\) 中的单项式理想，并假设 \(\{m_1, \ldots, m_k\}\) 是生成 \(I\) 的极小单项式集合，即每个 \(m_i\) 都是单项式，且 \(\{m_1, \ldots, m_k\}\) 的任何真子集都不生成 \(I\). 证明这些 \(m_i\)（\(1 \leq i \leq k\)）是唯一的。〔利用习题 10。〕

\item 在 \(R = F[x_1, \ldots, x_n]\) 上固定一个单项序。

（a）证明 \(\{g_1, \ldots, g_m\}\) 是 \(R\) 中理想 \(I\) 的极小 Gröbner 基当且仅当 \(\{\operatorname{LT}(g_1), \ldots, \operatorname{LT}(g_m)\}\) 是 \(\operatorname{LT}(I)\) 的极小生成集。

（b）证明 \(I\) 的极小 Gröbner 基的首项是唯一确定的，且 \(I\) 的任意两个极小 Gröbner 基中元素个数相同。〔利用（a）与上一习题。〕

\item 在 \(F[x_1, \ldots, x_n]\) 上固定一个单项序，并假设 \(G = \{g_1, \ldots, g_m\}\) 是非零理想 \(I\) 的一组生成元。证明若 \(S(g_i, g_j) \not\equiv 0 \bmod G\)，则理想 \((\operatorname{LT}(g_1), \ldots, \operatorname{LT}(g_m), \operatorname{LT}(S(g_i, g_j)))\) 严格大于理想 \((\operatorname{LT}(g_1), \ldots, \operatorname{LT}(g_m))\). 由此断定命题 \ref{pro:9.26} 之后描述的计算 Gröbner 基的算法在有限步内终止。〔利用习题 1。〕

\item 固定 \(F[x, y]\) 上的字典序 \(x > y\). 用 Buchberger 判别法证明 \(\{x^2y-y^2,\ x^3-xy\}\) 是理想 \(I = (x^2y-y^2,\ x^3-xy)\) 的 Gröbner 基。

\item 证明 \(\{x-y^3,\ y^5-y^6\}\) 是理想 \(I = (x-y^3,\ -x^2+xy^2)\) 关于 \(F[x, y]\) 中由 \(x > y\) 定义的字典序的约化 Gröbner 基。

\item 固定 \(F[x, y]\) 上的字典序 \(x > y\).

（a）证明 \(\{x^3-y,\ x^2y-y^2,\ xy^2-y^2,\ y^3-y^2\}\) 是理想 \(I = (-x^3+y,\ x^2y-y^2)\) 的约化 Gröbner 基。

（b）判断多项式 \(f = x^6-x^5y\) 是否为理想 \(I\) 的元素。

\item 固定 \(F[x, y, z]\) 上的字典序 \(x > y > z\). 证明 \(\{x^2+xy+z,\ xyz+z^2,\ xz^2,\ z^3\}\) 是理想 \(I = (x^2+xy+z,\ xyz+z^2)\) 的约化 Gröbner 基，并特别地断定首项理想 \(\operatorname{LT}(I)\) 需要四个生成元。

\item 固定 \(F[x, y]\) 上的字典序 \(x > y\). 用 Buchberger 判别法证明 \(\{x^2y-y^2,\ x^3-xy\}\) 是理想 \(I = (x^2y-y^2,\ x^3-xy)\) 的 Gröbner 基。

\item 设 \(I = (x^2-y,\ x^2y-z) \subset F[x, y, z]\).

（a）证明 \(\{x^2-y,\ y^2-z\}\) 是 \(I\) 关于由 \(x > y > z\) 定义的字典序的约化 Gröbner 基。

（b）证明 \(\{x^2-y,\ z-y^2\}\) 是 \(I\) 关于由 \(z > x > y\) 定义的字典序的约化 Gröbner 基（注意这些本质上与（a）中是同样的多项式）。

（c）证明 \(\{y-x^2,\ z-x^4\}\) 是 \(I\) 关于由 \(z > y > x\) 定义的字典序的约化 Gröbner 基。

\item 证明 \(F[x, y]\) 中的理想 \(I = (x^2y+xy^2-2y,\ x^2+xy-x+y^2-2y,\ xy^2-x-y+y^3)\) 与 \(J = (x-y^2,\ xy-y,\ x^2-y)\) 相等。

\item 用约化 Gröbner 基证明 \(F[x, y, z]\) 中的理想 \(I = (x^3-yz,\ yz+y)\) 与理想 \(J = (x^3z+x^3,\ x^3+y)\) 相等。

\item 证明理想 \(I = (x^2+xy^2,\ x^2-y^3,\ y^3-y^2)\) 关于字典序 \(x > y\) 的约化 Gröbner 基是 \(\{x^2-y^2,\ y^3-y^2,\ xy^2+y^2\}\).

\item 证明理想 \(I = (xy+y^2,\ x^2y+xy^2+x^2)\) 关于字典序 \(x > y\) 的约化 Gröbner 基是 \(\{x^2,\ xy+y^2,\ y^3\}\)，而关于字典序 \(y > x\) 的约化 Gröbner 基是 \(\{y^2+yx,\ x^2\}\).

\item 证明 \(\{x^3-y^3,\ x^2+xy^2+y^4,\ x^2y+xy^3+y^2\}\) 是推论 \ref{cor:9.28} 之后的例子中理想 \(I\) 关于 grlex 单项序的约化 Gröbner 基。（注意它给出 \(I\) 的三个生成元，而按字典序的例子中是两个，但次数更小。）

\item 设 \(I = (x^4-y^4+z^3-1,\ x^3+y^2+z^2-1)\). 证明 \(I\) 关于字典序 \(x > y > z\) 的约化 Gröbner 基有五个元素（这五个生成元中最大次数为 12，单项式项最多者为 35 项），关于字典序 \(y > z > x\) 有两个元素（最大次数为 6，项数最多为 8），而关于 grevlex 单项序的约化 Gröbner 基是 \(\{x^3+y^2+z^2-1,\ xy^2+xz^2-x+y^4-z^3+1\}\).

\item 在 \(\mathbb{C}\) 上解方程组 \(x^2-yz = 3\)、\(y^2-xz = 4\)、\(z^2-xy = 5\).

\item 求理想 \(I = (x^2+xy+y^2-1,\ x^2+4y^2-4)\) 关于字典序 \(x > y\) 的 Gröbner 基，并用它求出椭圆 \(x^2+xy+y^2 = 1\) 与椭圆 \(x^2+4y^2 = 4\) 在 \(\mathbb{R}^2\) 中的四个交点。

\item 用 Gröbner 基求出方程组 \(2x^3+2x^2y^2+3y^3 = 0\) 与 \(3x^5+2x^3y^3+2y^5 = 0\) 在 \(\mathbb{C}\) 上的全部六个解。

\item 用 Gröbner 基证明 \(F[x, y, z]\) 中 \((x, z) \cap (y^2, x-yz) = (xy, x-yz)\).

\item 用 Gröbner 基计算 \(F[x, y]\) 中理想 \((x^3y-xy^2+1,\ x^2y^2-y^3-1)\) 与 \((x^2-y^2,\ x^3+y^3)\) 的交。
\end{enumerate}

\noindent 下面四个习题讨论环 \(R\) 中两个理想 \(I\) 与 \(J\) 的\textbf{理想商}。

\begin{definition}
环 \(R\) 中两个理想 \(I, J\) 的\textbf{理想商} \((I : J)\) 是理想
\[
(I : J) = \{r \in R \mid rJ \subseteq I\}.
\]
\end{definition}

\begin{enumerate}
\setcounter{enumi}{33}
\item （a）假设 \(R\) 是整环，\(0 \neq f \in R\)，\(I\) 是 \(R\) 中的理想。证明若 \(\{g_1, \ldots, g_s\}\) 是理想 \(I \cap (f)\) 的生成元，则 \(\{g_1/f, \ldots, g_s/f\}\) 是理想商 \((I : (f))\) 的生成元。

（b）若 \(I\) 是交换环 \(R\) 中的理想，\(f_1, \ldots, f_s \in R\)，证明理想商 \((I : (f_1, \ldots, f_s))\) 是理想 \(\bigcap_{i=1}^s (I : (f_i))\).

\item 若 \(I = (x^2y+z^3,\ x+y^3-z,\ 2y^4z-yz^2-z^3)\)、\(J = (x^2y^5,\ x^3z^4,\ y^3z^7)\) 是 \(\mathbb{Q}[x, y, z]\) 中的理想，证明 \((I : J)\) 是理想 \((z^2,\ y+x-z)\).〔利用上一习题与命题 \ref{pro:9.30}。〕

\item 假设 \(K\) 是 \(R\) 中的理想，\(I\) 是包含 \(K\) 的理想，\(J\) 是任意理想。若 \(\bar{I}\) 和 \(\bar{J}\) 表示 \(I\) 和 \(J\) 在商环 \(R/K\) 中的像，证明 \(\overline{(I : J)} = (\bar{I} : \bar{J})\)，其中 \(\overline{(I : J)}\) 是理想商 \((I : J)\) 在 \(R/K\) 中的像。

\item 设 \(K\) 是 \(R = \mathbb{Q}[y, z]\) 中的理想 \((y^5-z^4)\). 对下列每一对理想 \(I\) 与 \(J\)，用前两个习题连同命题 \ref{pro:9.30} 验证环 \(R/K\) 中的理想商 \((\bar{I} : \bar{J})\)：

（i）\(I = (y^3,\ y^5-z^4)\)，\(J = (z)\)，\((\bar{I} : \bar{J}) = (\bar{y}^3, \bar{z})\)；

（ii）\(I = (y^3,\ z^5-y^5-z^4)\)，\(J = (z)\)，\((\bar{I} : \bar{J}) = (\bar{y}^3, \bar{z}^3)\)；

（iii）\(I = (y,\ y^3,\ y^5-z^4)\)，\(J = (1)\)，\((\bar{I} : \bar{J}) = (\bar{y}, \bar{z})\).
\end{enumerate}

\noindent 习题 38 到 44 发展 \(F[x_1, \ldots, x_n]\) 中单项式理想的一些补充的初等性质。由希尔伯特基定理可知理想都是有限生成的，然而在这些习题中不必假设这一点——对有限生成或无限生成的理想论证是相同的。这些习题可以用来给出希尔伯特基定理的一个独立证明（习题 44）。在这些习题中，\(M\) 和 \(N\) 是单项式理想，其单项式生成元分别为 \(\{m_i \mid i \in I\}\) 和 \(\{n_j \mid j \in J\}\)（对某些指标集 \(I\) 和 \(J\)）。

\begin{enumerate}
\setcounter{enumi}{37}
\item 通过证明 \(M+N = (m_i, n_j \mid i \in I, j \in J)\) 和 \(MN = (m_in_j \mid i \in I, j \in J)\)，证明两个单项式理想的和与积都是单项式理想。

\item 证明若 \(\{M_s \mid s \in S\}\) 是任意非空的单项式理想族，且按包含关系全序，则 \(\bigcup_{s \in S} M_s\) 是单项式理想。（特别地，任何递增的单项式理想序列的并都是单项式理想，参见第 7.3 节习题 19。）

\item 通过证明 \(M \cap N = (e_{i,j} \mid i \in I, j \in J)\) 证明两个单项式理想的交是单项式理想，其中 \(e_{i,j}\) 是 \(m_i\) 与 \(n_j\) 的最小公倍。〔利用习题 10。〕

\item 证明对任意单项式 \(n\)，理想商 \((M : (n))\) 是 \((m_i/d_i \mid i \in I)\)，其中 \(d_i\) 是 \(m_i\) 与 \(n\) 的最大公因子（参见习题 34）。证明若 \(N\) 是有限生成的，则两个单项式理想的理想商 \((M : N)\) 是单项式理想。

\item （a）证明 \(M\) 是单项式素理想当且仅当 \(M = (S)\) 对 \(\{x_1, x_2, \ldots, x_n\}\) 的某个子集 \(S\) 成立。（特别地，单项式素理想只有有限多个，且每个都是有限生成的。）

（b）证明 \((x_1, \ldots, x_n)\) 是唯一的单项式极大理想。

\item （\textbf{Dickson 引理——希尔伯特基定理的一个特殊情形}）按下述方式证明 \(F[x_1, \ldots, x_n]\) 中每个单项式理想都是有限生成的。设 \(S = \{N \mid N\) 是不是有限生成的单项式理想\(\}\)，并反设 \(S \neq \varnothing\).

（a）证明 \(S\) 含有一个极大元 \(M\).〔利用习题 39 与佐恩引理。〕

（b）证明存在不属于 \(M\) 的单项式 \(x\) 使得 \(xM \subseteq M\).〔利用习题 42（a）。〕

（c）对（b）中的 \(x\)，证明 \(M\) 含有一个有限生成的单项式理想 \(M_0\)，使得 \(M_0+(x) = M+(x)\) 且 \(M = M_0 + (x)(M : (x))\)，其中 \((M : (x))\) 是习题 41 中定义的（单项式）理想，而 \((x)(M : (x))\) 是这两个理想之积。由此断定 \(M\) 是有限生成的，这个矛盾证明 \(S = \varnothing\).〔利用 \(M\) 的极大性与前面的习题。〕

\item 若 \(I\) 是 \(F[x_1, \ldots, x_n]\) 中的非零理想，用 Dickson 引理证明 \(\operatorname{LT}(I)\) 是有限生成的。由此断定 \(I\) 有 Gröbner 基，并推出希尔伯特基定理。〔参见命题 \ref{pro:9.24}。〕

\item （\textbf{图的 \(n\)-着色}）大小为 \(N\) 的有限图 \(\mathcal{G}\) 是顶点集 \(i \in \{1, 2, \ldots, N\}\) 与连接顶点 \(i\) 与顶点 \(j\) 的边 \((i, j)\) 的集合。\(\mathcal{G}\) 的一个 \textbf{\(n\)-着色}是把 \(n\) 种颜色之一分配给每个顶点，使得由一条边相连的顶点颜色不同。设 \(F\) 是任意含有至少 \(n\) 个元素的域。若我们对每个顶点 \(i\) 引入一个变量 \(x_i\)，并通过从 \(F\) 中选取 \(n\) 个互异元素组成的集合 \(S\) 来表示这 \(n\) 种颜色，则 \(\mathcal{G}\) 的一个 \(n\)-着色等价于对每个 \(i = 1, 2, \ldots, N\) 指定一个值 \(x_i = a_i\)，其中 \(a_i \in S\)，且若 \((i, j)\) 是 \(\mathcal{G}\) 中的边则 \(a_i \neq a_j\). 若 \(f(x) = \prod_{a \in S}(x-a)\) 是 \(F[x]\) 中以 \(S\) 中元素为根的 \(n\) 次多项式，则 \(x_i = a_i\)（某个 \(a_i \in S\)）等价于 \(x_i\) 是方程 \(f(x_i) = 0\) 的解。于是 \(a_i \neq a_j\) 就是 \(f(x_i) = f(x_j)\) 而 \(x_i \neq x_j\)，即 \(x_i\) 与 \(x_j\) 满足方程 \(g(x_i, x_j) = 0\)，其中 \(g(x_i, x_j)\) 是 \(F[x_i, x_j]\) 中的多项式 \((f(x_i)-f(x_j))/(x_i-x_j)\). 由此，求 \(\mathcal{G}\) 的一个 \(n\)-着色等价于解方程组
\[
\begin{cases}
f(x_i) = 0, & i = 1, 2, \ldots, N, \\
g(x_i, x_j) = 0, & \mathcal{G} \text{ 的所有边 } (i, j)
\end{cases}
\]
（还要注意只要 \(g(a_i, a_j) = 0\) 推出 \(a_i \neq a_j\)，我们可以使用任何多项式 \(g\)）。由“希尔伯特零点定理”（参见第 15.3 节推论 33）可知，这个方程组有解——从而 \(\mathcal{G}\) 有 \(n\)-着色——除非 \(F[x_1, x_2, \ldots, x_N]\) 中由多项式 \(f(x_i)\)（\(i = 1, 2, \ldots, N\)）连同图 \(\mathcal{G}\) 中所有边 \((i, j)\) 的多项式 \(g(x_i, x_j)\) 生成的理想 \(I\) 不是真理想。这等价于说 \(I\) 的约化 Gröbner 基（关于任何单项序）就是 \(\{1\}\). 进一步，当 \(n\)-着色存在时，像命题 \ref{pro:9.29} 后的例子那样解这个方程组就给出 \(\mathcal{G}\) 的一个显式着色。

域 \(F\) 与集合 \(S\) 有许多可能的选择。例如，取任意含有 \(n\) 个互异的 \(n\) 次单位根的集合 \(S\) 的域 \(F\)，此时 \(f(x) = x^n-1\)，我们可以取 \(g(x_i, x_j) = (x_i^n-x_j^n)/(x_i-x_j) = x_i^{n-1}+x_i^{n-2}x_j+\cdots+x_ix_j^{n-2}+x_j^{n-1}\)；或者取 \(F = \mathbb{F}_p\)（\(p \geq n\) 是素数）的任意子集 \(S\)（特别地，当 \(n = p\) 时由费马小定理有 \(f(x) = x^p-x\) 且 \(g(x_i, x_j) = (x_i-x_j)^{p-1}-1\)）。

（a）考虑有八个顶点、14 条边 \((1,3), (1,4), (1,5), (2,4), (2,7), (2,8), (3,4), (3,6), (3,8), (4,5), (5,6), (6,7), (6,8), (7,8)\) 的图 \(\mathcal{G}\) 的一个可能的 3-着色。取 \(F = \mathbb{F}_3\)，“颜色”为 \(0, 1, 2 \in \mathbb{F}_3\)，并设顶点 1 被染成 0. 此时 \(f(x) = x(x-1)(x-2) = x^3-x \in \mathbb{F}_3[x]\)，而 \(g(x_i, x_j) = x_i^2+x_ix_j+x_j^2-1\). 若 \(I\) 是由 \(x_1\)、\(x_i^3-x_i\)（\(2 \leq i \leq 8\)）以及 \(\mathcal{G}\) 中所有边 \((i, j)\) 的多项式 \(g(x_i, x_j)\) 生成的理想，证明 \(I\) 关于字典单项序 \(x_1 > x_2 > \cdots > x_8\) 的约化 Gröbner 基是 \(\{x_1, x_2, x_3+x_8, x_4+2x_8, x_5+x_8, x_6, x_7+x_8, x_8^2+2\}\). 由此断定 \(\mathcal{G}\) 有两种不同的 3-着色，它们由顶点 8 的着色决定（顶点 8 必须被染成 \(\mathbb{F}_3\) 中的非零元素），并具体给出 \(\mathcal{G}\) 的这些着色。证明若把边 \((3,7)\) 加到 \(\mathcal{G}\) 上，则该图不能被 3-着色。

（b）取 \(F = \mathbb{F}_5\)，“颜色”为 \(1, 2, 3, 4 \in \mathbb{F}_5\)，于是 \(f(x) = x^4-1\)，我们可以取 \(g(x_i, x_j) = x_i^3+x_i^2x_j+x_ix_j^2+x_j^3\). 证明有五个顶点、9 条边 \((1,3), (1,4), (1,5), (2,3), (2,4), (2,5), (3,4), (3,5), (4,5)\) 的图 \(\mathcal{G}\)（“五个顶点的完全图去掉一条边”）可以被 4-着色，但不能被 3-着色。

（c）用 Gröbner 基证明有九个顶点、22 条边 \((1,4), (1,6), (1,7), (1,8), (2,3), (2,4), (2,6), (2,7), (3,5), (3,7), (3,9), (4,5), (4,6), (4,7), (4,9), (5,6), (5,7), (5,8), (5,9), (6,7), (6,9), (7,8)\) 的图 \(\mathcal{G}\) 在颜色的置换意义下恰有四种 4-着色（故共有 96 个 4-着色）。证明若把边 \((1,5)\) 加上去，则 \(\mathcal{G}\) 不能被 4-着色。
\end{enumerate}

